% Anomaly-completed topological holography: extensions
% Full Morita theory, gravitational regime, critical points,
% symplectic polarization of complementarity

\providecommand{\Gg}{\mathfrak{G}}
\providecommand{\Sgmod}{S_g}
\providecommand{\Mbar}{\overline{\mathcal{M}}}
\providecommand{\ChirHoch}{\operatorname{ChirHoch}}

\subsection{Full Morita theory of the genus-Clifford completion}
% label removed: subsec:tholog-full-morita-theory

\index{Morita equivalence!genus-Clifford|textbf}
\index{spinor module|textbf}
\index{genus-Clifford completion!Morita theory}

Theorem~\ref{thm:tholog-morita-triviality} established that
$\Gg_g(B_\Theta)[u^{-1}] \cong \operatorname{Mat}_{2^g}(R)$
with $R = B_\Theta[u^{-1}]$.
The proof used the tensor-product splitting of the genus-$g$
algebra into $g$ copies of the genus-$1$ algebra, each
separately identified with $\operatorname{Mat}_2(R)$.
The present section constructs the Morita equivalence
explicitly at each genus, building the spinor module $\Sgmod$,
verifying the endomorphism computation, and spelling out
both functors of the equivalence.

\subsubsection{The Clifford algebra over a central base}
% label removed: subsubsec:tholog-clifford-over-base

Fix a unital associative dg algebra $R$ and a central
closed degree-$2$ element $u \in Z^2(R) \cap Z(R)$.
For $g \geq 1$, define
\[
\operatorname{Cl}_{2g}(R,u)
\;:=\;
\frac{R\langle \alpha_1,\beta_1,\dots,\alpha_g,\beta_g \rangle}
{\bigl\langle
  \alpha_i\alpha_j + \alpha_j\alpha_i,\;
  \beta_i\beta_j + \beta_j\beta_i,\;
  \alpha_i\beta_j + \beta_j\alpha_i - \delta_{ij}u
\bigr\rangle},
\]
with $|\alpha_i| = |\beta_i| = 1$ and $d\alpha_i = d\beta_i = 0$.
By definition, $\Gg_g(B_\Theta) = \operatorname{Cl}_{2g}(B_\Theta, u)$
with $u = \eta^2$.

\begin{lemma}[Poincar\'{e}--Birkhoff--Witt basis; \ClaimStatusProvedHere]
% label removed: lem:tholog-clifford-pbw
As a left $R$-module, $\operatorname{Cl}_{2g}(R,u)$ is free of rank
$4^g$ with basis
\[
\big\{
  \alpha_1^{a_1}\beta_1^{b_1}\cdots\alpha_g^{a_g}\beta_g^{b_g}
  \;\big|\;
  a_i, b_i \in \{0,1\}
\big\}.
\]
In particular, $\operatorname{Cl}_{2g}(R,u) \cong R^{4^g}$ as a
graded left $R$-module.
\end{lemma}

\begin{proof}
The Clifford relations reduce every monomial in the generators to
one in which each $\alpha_i$ and each $\beta_i$ appears at most
once (using $\alpha_i^2 = 0$, $\beta_i^2 = 0$), and the
generators for different handle indices graded-commute (the
cross-handle anticommutators are all zero).
The ordered monomials span $\operatorname{Cl}_{2g}(R,u)$
over~$R$.
Freeness follows because the Clifford algebra
$\operatorname{Cl}_{2g}(R,u)$ is filtered by $\eta$-degree,
with the filtration $F^p := \sum_{k \leq p} R \cdot
\{\text{monomials of total generator degree} \leq k\}$.
The associated graded is isomorphic to the exterior algebra
$\Lambda^{2g}(\alpha_1,\beta_1,\dots,\alpha_g,\beta_g)$
over~$R$ (since the Clifford relations become anticommutation
relations modulo lower filtration), which is free of rank~$4^g$
over~$R$.  The filtration is exhaustive (every element has
finite generator degree) and separated (the intersection
$\bigcap_p F^p = F^0 = R$ contributes only the unit),
so $\operatorname{Cl}_{2g}(R,u)$ inherits freeness of
rank~$4^g$ from its associated graded.
\end{proof}

\subsubsection{The spinor module: genus $g = 1$}
% label removed: subsubsec:tholog-spinor-g1

\begin{definition}[Spinor module at genus 1]
% label removed: def:tholog-spinor-g1
\index{spinor module!genus 1}
Let $R = B_\Theta[u^{-1}]$.
The \emph{genus-$1$ spinor module} is
\[
S_1 \;:=\; R e_+ \oplus R e_-,
\]
where $|e_+| = 0$, $|e_-| = -1$, $de_+ = de_- = 0$, equipped with
the left $\Gg_1(R)$-action defined by
\begin{equation}% label removed: eq:spinor-g1-action
\alpha \cdot e_+ = u\,e_-,\qquad
\alpha \cdot e_- = 0,\qquad
\beta \cdot e_+ = 0,\qquad
\beta \cdot e_- = e_+.
\end{equation}
(Here and below we write $\alpha = \alpha_1$, $\beta = \beta_1$
for the single handle pair.)
\end{definition}

\begin{proposition}[Spinor module is well defined; \ClaimStatusProvedHere]
% label removed: prop:tholog-spinor-g1-well-defined
The formulas~\eqref{eq:spinor-g1-action} extend to a well-defined
left dg $\Gg_1(R)$-module structure on $S_1$.
\end{proposition}

\begin{proof}
We verify the three defining relations of $\Gg_1(R)$.

\emph{Relation $\alpha^2 = 0$.}
\begin{align*}
\alpha^2 \cdot e_+ &= \alpha(\alpha \cdot e_+) = \alpha(u\,e_-)
= u(\alpha \cdot e_-) = u \cdot 0 = 0, \\
\alpha^2 \cdot e_- &= \alpha(\alpha \cdot e_-) = \alpha(0) = 0.
\end{align*}
Since $\alpha^2 = 0$ on both generators, the relation holds.

\emph{Relation $\beta^2 = 0$.}
\begin{align*}
\beta^2 \cdot e_+ &= \beta(\beta \cdot e_+) = \beta(0) = 0, \\
\beta^2 \cdot e_- &= \beta(\beta \cdot e_-) = \beta(e_+) = 0.
\end{align*}
The last equality uses $\beta \cdot e_+ = 0$.

\emph{Relation $\alpha\beta + \beta\alpha = u$.}
On $e_+$:
\[
(\alpha\beta + \beta\alpha) \cdot e_+
= \alpha(\beta \cdot e_+) + \beta(\alpha \cdot e_+)
= \alpha(0) + \beta(u\,e_-)
= 0 + u(\beta \cdot e_-)
= u \cdot e_+.
\]
On $e_-$:
\[
(\alpha\beta + \beta\alpha) \cdot e_-
= \alpha(\beta \cdot e_-) + \beta(\alpha \cdot e_-)
= \alpha(e_+) + \beta(0)
= u\,e_- + 0
= u \cdot e_-.
\]
In both cases, $(\alpha\beta + \beta\alpha)$ acts as
multiplication by $u$.

All three Clifford relations are satisfied.
It remains to verify that the differential is compatible with
the module structure.  Since $d(\alpha) = d(\beta) = d(e_\pm) = 0$
by hypothesis, the Leibniz rule
$d(a \cdot m) = (da) \cdot m + (-1)^{|a|} a \cdot (dm)$
is trivially satisfied: both terms on the right vanish for
every generator pair $(a, m)$.
\end{proof}

\begin{proposition}[Endomorphism computation at genus 1; \ClaimStatusProvedHere]
% label removed: prop:tholog-end-g1
There is an isomorphism of dg algebras
\[
\underline{\operatorname{End}}_R(S_1)
\;\cong\;
\operatorname{Mat}_{2 \times 2}(R).
\]
The algebra map
$\rho : \Gg_1(R) \to \underline{\operatorname{End}}_R(S_1)$
induced by the $\Gg_1(R)$-action on $S_1$ is an isomorphism.
\end{proposition}

\begin{proof}
Since $S_1 = R e_+ \oplus R e_-$ is a free $R$-module of
rank~$2$, the endomorphism algebra is
$\underline{\operatorname{End}}_R(S_1) \cong \operatorname{Mat}_2(R)$
as graded $R$-algebras.
Write the elementary matrix $E_{ij}$ for the endomorphism
$e_j \mapsto e_i$, $e_k \mapsto 0$ for $k \neq j$
(with $e_1 = e_+$, $e_2 = e_-$).

Under $\rho$, the action formulas give:
\begin{alignat*}{2}
\rho(\alpha) &: e_+ \mapsto u\,e_-,\; e_- \mapsto 0
&\qquad &= u\,E_{21}, \\
\rho(\beta) &: e_+ \mapsto 0,\; e_- \mapsto e_+
&\qquad &= E_{12}, \\
\rho(\alpha\beta) &: e_+ \mapsto 0,\; e_- \mapsto u\,e_-
&\qquad &= u\,E_{22}.
\end{alignat*}
Compute $\rho(\alpha\beta)$ directly:
$\rho(\alpha\beta) \cdot e_+ = \alpha(\beta \cdot e_+)
= \alpha(0) = 0$, and
$\rho(\alpha\beta) \cdot e_- = \alpha(\beta \cdot e_-)
= \alpha(e_+) = u\,e_-$.
Also, $\rho(1) = E_{11} + E_{22}$.

The four elements $\{1, \alpha, \beta, \alpha\beta\}$ form
a basis of $\Gg_1(R)$ over $R$
(Lemma~\ref{lem:tholog-clifford-pbw}).
Their images are
\[
E_{11} + E_{22},\quad
u\,E_{21},\quad
E_{12},\quad
u\,E_{22}.
\]
Since $u$ is invertible in $R$, we recover
$E_{21} = u^{-1}\rho(\alpha)$,
$E_{12} = \rho(\beta)$,
$E_{22} = u^{-1}\rho(\alpha\beta)$, and
$E_{11} = \rho(1) - u^{-1}\rho(\alpha\beta)$.
Thus $\rho$ is surjective.
As both sides are free of rank~$4$ over $R$, surjectivity
implies $\rho$ is an isomorphism.
\end{proof}

\begin{corollary}[Morita bimodule at genus 1; \ClaimStatusProvedHere]
% label removed: cor:tholog-morita-bimodule-g1
The spinor module $S_1$ is a $(\Gg_1(R), R)$-bimodule that
implements the Morita equivalence
\[
D\!\bigl(\Gg_1(R)\text{-}\mathbf{mod}\bigr)
\;\simeq\;
D\!\bigl(R\text{-}\mathbf{mod}\bigr).
\]
The equivalence functor is $M \mapsto S_1 \otimes_R M$ with
inverse $N \mapsto S_1^\vee \otimes_{\Gg_1(R)} N$, where
$S_1^\vee = \underline{\operatorname{Hom}}_R(S_1, R)$.
\end{corollary}

\begin{proof}
Proposition~\ref{prop:tholog-end-g1} identifies
$\Gg_1(R) \cong \operatorname{End}_R(S_1)$.
Since $S_1$ is a finitely generated projective (in fact free)
$R$-module of rank $2$, the standard Morita theorem
(Lam, \emph{Lectures on Modules and Rings}, Theorem~18.30)
gives the derived equivalence.
The inverse functor is
$N \mapsto S_1^\vee \otimes_{\Gg_1(R)} N
= \operatorname{Hom}_R(S_1, R) \otimes_{\Gg_1(R)} N$,
and the unit and counit of the adjunction are the standard
evaluation and coevaluation maps.
\end{proof}

\subsubsection{The spinor module: genus $g = 2$}
% label removed: subsubsec:tholog-spinor-g2

\begin{definition}[Spinor module at genus 2]
% label removed: def:tholog-spinor-g2
\index{spinor module!genus 2}
Define
\[
S_2 \;:=\; S_1 \otimes_R S_1'
\;=\;
R^4,
\]
where $S_1 = R e_+ \oplus R e_-$ carries the handle pair
$(\alpha_1, \beta_1)$ and
$S_1' = R e_+' \oplus R e_-'$ carries a second copy
$(\alpha_2, \beta_2)$ with the same formulas as~\eqref{eq:spinor-g1-action}.
The graded tensor product gives a free $R$-module of rank $4$
with basis
\[
e_{++} := e_+ \otimes e_+',\quad
e_{+-} := e_+ \otimes e_-',\quad
e_{-+} := e_- \otimes e_+',\quad
e_{--} := e_- \otimes e_-'.
\]
The degrees are
$|e_{++}| = 0$, $|e_{+-}| = -1$, $|e_{-+}| = -1$, $|e_{--}| = -2$.
\end{definition}

\begin{proposition}[Genus-2 Clifford relations; \ClaimStatusProvedHere]
% label removed: prop:tholog-spinor-g2-relations
The actions
\begin{alignat*}{2}
\alpha_1 &= u\,E_{21} \otimes \operatorname{id},&\qquad
\beta_1 &= E_{12} \otimes \operatorname{id}, \\
\alpha_2 &= \varepsilon \otimes u\,E_{21},&\qquad
\beta_2 &= \varepsilon \otimes E_{12},
\end{alignat*}
where $\varepsilon = (-1)^{|\cdot|}$ is the grading operator
on $S_1$ (acting as $+1$ on $e_+$ and $-1$ on $e_-$),
define a well-defined left $\Gg_2(R)$-module structure on~$S_2$.
\end{proposition}

\begin{proof}
The grading operator $\varepsilon$ satisfies
$\varepsilon^2 = \operatorname{id}$ and anticommutes with both
$E_{21}$ and $E_{12}$ (since these shift degree by $\pm 1$):
\[
\varepsilon E_{21} = -E_{21} \varepsilon,\qquad
\varepsilon E_{12} = -E_{12} \varepsilon.
\]

\emph{Same-handle relations.}
For handle $1$: the operators $\alpha_1 = u\,E_{21}\otimes\operatorname{id}$
and $\beta_1 = E_{12}\otimes\operatorname{id}$ satisfy
$\alpha_1^2 = u^2(E_{21})^2\otimes\operatorname{id} = 0$,
$\beta_1^2 = (E_{12})^2\otimes\operatorname{id} = 0$, and
\begin{align*}
\alpha_1\beta_1 + \beta_1\alpha_1
&= u(E_{21}E_{12} + E_{12}E_{21})\otimes\operatorname{id} \\
&= u\,(E_{22} + E_{11})\otimes\operatorname{id}
= u\,\operatorname{id}_{S_1}\otimes\operatorname{id}_{S_1'}
= u\,\operatorname{id}_{S_2}.
\end{align*}
For handle $2$: the operators
$\alpha_2 = \varepsilon \otimes u\,E_{21}$
and $\beta_2 = \varepsilon \otimes E_{12}$ satisfy
$\alpha_2^2 = \varepsilon^2 \otimes u^2(E_{21})^2 = \operatorname{id} \otimes 0 = 0$,
$\beta_2^2 = \varepsilon^2 \otimes (E_{12})^2 = 0$, and
\begin{align*}
\alpha_2\beta_2 + \beta_2\alpha_2
&= \varepsilon^2 \otimes u(E_{21}E_{12} + E_{12}E_{21}) \\
&= \operatorname{id} \otimes u\,\operatorname{id}
= u\,\operatorname{id}_{S_2}.
\end{align*}

\emph{Cross-handle relations.}
We must verify $\alpha_1\alpha_2 + \alpha_2\alpha_1 = 0$,
$\beta_1\beta_2 + \beta_2\beta_1 = 0$,
$\alpha_1\beta_2 + \beta_2\alpha_1 = 0$,
$\alpha_2\beta_1 + \beta_1\alpha_2 = 0$.

For $\alpha_1\alpha_2 + \alpha_2\alpha_1$:
\begin{align*}
&(u\,E_{21}\otimes\operatorname{id})
(\varepsilon\otimes u\,E_{21})
+
(\varepsilon\otimes u\,E_{21})
(u\,E_{21}\otimes\operatorname{id}) \\
&\quad=\;
u(E_{21}\varepsilon)\otimes(u\,E_{21})
+
u(\varepsilon E_{21})\otimes(u\,E_{21}) \\
&\quad=\;
u^2(E_{21}\varepsilon + \varepsilon E_{21})\otimes E_{21}
\;=\; 0,
\end{align*}
since $E_{21}\varepsilon + \varepsilon E_{21} = 0$
by the anticommutation above.
The remaining three cross-handle computations are identical
in structure, each reducing to
$E_{ij}\varepsilon + \varepsilon E_{ij} = 0$
for an off-diagonal $E_{ij}$.
\end{proof}

\begin{proposition}[Endomorphism computation at genus 2; \ClaimStatusProvedHere]
% label removed: prop:tholog-end-g2
There is an isomorphism of dg algebras
\[
\Gg_2(R)
\;\cong\;
\underline{\operatorname{End}}_R(S_2)
\;\cong\;
\operatorname{Mat}_{4 \times 4}(R).
\]
\end{proposition}

\begin{proof}
By Lemma~\ref{lem:tholog-clifford-pbw}, $\Gg_2(R)$ is free of
rank $4^2 = 16$ over $R$.
Likewise $\underline{\operatorname{End}}_R(S_2) \cong \operatorname{Mat}_4(R)$
is free of rank $16$.
The algebra map
$\rho : \Gg_2(R) \to \operatorname{Mat}_4(R)$
is the tensor product of the two genus-$1$ representations.
By Proposition~\ref{prop:tholog-end-g1}, each factor is an
isomorphism after inverting $u$.
The tensor product of two matrix-algebra isomorphisms gives
\[
\operatorname{Mat}_2(R) \otimes_R \operatorname{Mat}_2(R)
\;\cong\;
\operatorname{Mat}_4(R),
\]
and the composite
$\Gg_2(R) = \Gg_1(R) \otimes_R \Gg_1(R)
\xrightarrow{\rho_1 \otimes \rho_2}
\operatorname{Mat}_2(R) \otimes_R \operatorname{Mat}_2(R)
\cong \operatorname{Mat}_4(R)$
is an isomorphism.
\end{proof}

We record the explicit action on the four basis vectors
for later use.

\begin{computation}[Explicit genus-2 spinor action; \ClaimStatusProvedHere]
% label removed: comp:tholog-spinor-g2-explicit
\index{spinor module!genus 2!explicit action}
In the ordered basis $(e_{++}, e_{+-}, e_{-+}, e_{--})$:
\[
\rho(\alpha_1) = \begin{pmatrix}
0 & 0 & 0 & 0 \\
0 & 0 & 0 & 0 \\
u & 0 & 0 & 0 \\
0 & u & 0 & 0
\end{pmatrix},\qquad
\rho(\beta_1) = \begin{pmatrix}
0 & 0 & 1 & 0 \\
0 & 0 & 0 & 1 \\
0 & 0 & 0 & 0 \\
0 & 0 & 0 & 0
\end{pmatrix},
\]
\[
\rho(\alpha_2) = \begin{pmatrix}
0 & 0 & 0 & 0 \\
u & 0 & 0 & 0 \\
0 & 0 & 0 & 0 \\
0 & 0 & -u & 0
\end{pmatrix},\qquad
\rho(\beta_2) = \begin{pmatrix}
0 & 1 & 0 & 0 \\
0 & 0 & 0 & 0 \\
0 & 0 & 0 & -1 \\
0 & 0 & 0 & 0
\end{pmatrix}.
\]
Direct multiplication confirms all six Clifford relations.
For instance,
\[
\rho(\alpha_1)\rho(\beta_2) + \rho(\beta_2)\rho(\alpha_1) = 0
\]
because the $(3,2)$-entry of the first product is $u$ while
the $(3,2)$-entry of the second is $-u$, and all other
nonzero entries cancel similarly.
\end{computation}

\subsubsection{The spinor module: genus $g = 3$}
% label removed: subsubsec:tholog-spinor-g3

\begin{definition}[Spinor module at genus 3]
% label removed: def:tholog-spinor-g3
\index{spinor module!genus 3}
Define
\[
S_3 \;:=\; S_2 \otimes_R S_1''
\;=\; R^8,
\]
where $S_1'' = R e_+'' \oplus R e_-''$ carries the third handle
pair $(\alpha_3, \beta_3)$.
The basis vectors are
$e_{\epsilon_1\epsilon_2\epsilon_3}$ with
$\epsilon_i \in \{+,-\}$, and the $\Gg_3(R)$-action is the
graded tensor product of the three genus-$1$ representations:
handle $i$ acts on the $i$-th tensor factor, inserted with
grading operators on the preceding factors.
\end{definition}

\begin{proposition}[Endomorphism computation at genus 3; \ClaimStatusProvedHere]
% label removed: prop:tholog-end-g3
There is an isomorphism of dg algebras
\[
\Gg_3(R)
\;\cong\;
\underline{\operatorname{End}}_R(S_3)
\;\cong\;
\operatorname{Mat}_{8 \times 8}(R).
\]
\end{proposition}

\begin{proof}
Iterate the tensor-product argument of
Proposition~\ref{prop:tholog-end-g2}: three copies of
$\operatorname{Mat}_2(R)$ tensor to $\operatorname{Mat}_8(R)$,
and the composite map is an isomorphism by the
$g$-fold product of the genus-$1$ isomorphisms.
\end{proof}

\subsubsection{The general spinor module and the full Morita equivalence}
% label removed: subsubsec:tholog-general-spinor

\begin{definition}[General spinor module]
% label removed: def:tholog-spinor-general
\index{spinor module!general genus}
For $g \geq 1$, the \emph{genus-$g$ spinor module} is
\[
S_g \;:=\; S_1^{\otimes_R g}
\;=\;
R^{2^g},
\]
with basis $\{e_{\epsilon_1 \cdots \epsilon_g}\}$
indexed by $(\epsilon_1,\dots,\epsilon_g) \in \{+,-\}^g$.
The $\Gg_g(R)$-action on the $i$-th handle pair
$(\alpha_i, \beta_i)$ acts on the $i$-th tensor factor,
with the grading operator $\varepsilon_j = (-1)^{|\cdot|}$
on each preceding factor $j < i$:
\begin{align}
\alpha_i &= \varepsilon_1 \otimes \cdots \otimes \varepsilon_{i-1}
\otimes u\,E_{21} \otimes \operatorname{id} \otimes \cdots \otimes
\operatorname{id},
% label removed: eq:general-spinor-alpha \\
\beta_i &= \varepsilon_1 \otimes \cdots \otimes \varepsilon_{i-1}
\otimes E_{12} \otimes \operatorname{id} \otimes \cdots \otimes
\operatorname{id}.
% label removed: eq:general-spinor-beta
\end{align}
\end{definition}

\begin{theorem}[Full Morita equivalence at genus $g$; \ClaimStatusProvedHere]
% label removed: thm:tholog-full-morita-g
For every $g \geq 1$ and $R = B_\Theta[u^{-1}]$:
\begin{enumerate}[label=\textup{(\roman*)}]
\item The algebra map
  $\rho_g : \Gg_g(R) \to \underline{\operatorname{End}}_R(S_g)
  \cong \operatorname{Mat}_{2^g}(R)$
  defined by~\eqref{eq:general-spinor-alpha}--\eqref{eq:general-spinor-beta}
  is an isomorphism of dg algebras.
\item The spinor module $S_g$ implements a Morita equivalence
  \[
  D\!\bigl(\Gg_g(R)\text{-}\mathbf{mod}\bigr)
  \;\simeq\;
  D\!\bigl(R\text{-}\mathbf{mod}\bigr).
  \]
\item The equivalence functors are:
  \[
  \Phi_g : M \longmapsto S_g \otimes_R M,\qquad
  \Psi_g : N \longmapsto S_g^\vee \otimes_{\Gg_g(R)} N.
  \]
\item The natural isomorphisms
  $\Psi_g \circ \Phi_g \cong \operatorname{id}$
  and $\Phi_g \circ \Psi_g \cong \operatorname{id}$
  are induced by the trace pairing
  $S_g^\vee \otimes_R S_g \cong R$
  and the coevaluation
  $R \to S_g \otimes_R S_g^\vee \cong \operatorname{End}_R(S_g)$.
\end{enumerate}
\end{theorem}

\begin{proof}
Part~(i).
The algebra $\Gg_g(R)$ splits as the tensor product of $g$
copies of $\Gg_1(R)$ (Definition~\ref{def:tholog-genus-clifford}:
generators for different handles graded-commute, and the
Clifford relation at each handle involves only the
local pair).
Proposition~\ref{prop:tholog-end-g1} gives
$\rho_1 : \Gg_1(R) \xrightarrow{\sim} \operatorname{Mat}_2(R)$.
The map $\rho_g = \rho_1^{\otimes g}$ is therefore an
isomorphism onto
$\operatorname{Mat}_2(R)^{\otimes g} \cong \operatorname{Mat}_{2^g}(R)$.

It remains to verify that the insertion of grading operators
$\varepsilon_j$ in~\eqref{eq:general-spinor-alpha}--\eqref{eq:general-spinor-beta}
is compatible with the tensor-product identification.
The grading operator on $S_1 = R e_+ \oplus R e_-$ acts as
$\varepsilon(e_+) = e_+$ and $\varepsilon(e_-) = -e_-$.
In the matrix basis this is
$\varepsilon = E_{11} - E_{22} = \begin{psmallmatrix}1&0\\0&-1\end{psmallmatrix}$.
The insertion
$\varepsilon_1 \otimes \cdots \otimes \varepsilon_{i-1}
\otimes E_{21} \otimes \operatorname{id} \otimes \cdots$
is the standard ``Jordan--Wigner transformation'' that
converts the graded tensor product of Clifford algebras into
a single matrix algebra.
It is classical that this map is an algebra homomorphism whose
image is all of $\operatorname{Mat}_{2^g}(R)$.

Part~(ii).
Since $S_g$ is a free $R$-module of finite rank $2^g$,
and $\Gg_g(R) \cong \operatorname{End}_R(S_g)$, the Morita
equivalence follows by the standard criterion:
a progenerator of rank $2^g$.

Part~(iii).
The equivalence functors are defined explicitly.
The forward functor is
$\Phi_g(M) := S_g \otimes_R M$,
which is a left $\Gg_g(R)$-module via the action on the
first tensor factor.
The inverse functor is
$\Psi_g(N) := \operatorname{Hom}_{\Gg_g(R)}(S_g, N)
\cong S_g^\vee \otimes_{\Gg_g(R)} N$,
where $S_g^\vee = \underline{\operatorname{Hom}}_R(S_g, R)$.

Part~(iv).
The trace pairing $\operatorname{tr} : S_g^\vee \otimes_R S_g \to R$
is the evaluation map $f \otimes v \mapsto f(v)$.
In the basis $\{e_{\epsilon_1\cdots\epsilon_g}\}$
with dual basis $\{e^{\epsilon_1\cdots\epsilon_g}\}$,
the trace is
$\operatorname{tr} = \sum_{\epsilon \in \{+,-\}^g}
e^{\epsilon_1\cdots\epsilon_g} \otimes e_{\epsilon_1\cdots\epsilon_g}
\;\mapsto\; 2^g \in R$.
The coevaluation
$\operatorname{coev} : R \to S_g \otimes_R S_g^\vee
\cong \operatorname{End}_R(S_g)$
sends $1 \mapsto \operatorname{id}_{S_g}
= \sum_\epsilon e_\epsilon \otimes e^\epsilon$.
The composite $\Psi_g \circ \Phi_g(M)
= S_g^\vee \otimes_{\Gg_g(R)} (S_g \otimes_R M)
\cong (S_g^\vee \otimes_{\Gg_g(R)} S_g) \otimes_R M
\cong R \otimes_R M \cong M$
via the trace pairing,
and $\Phi_g \circ \Psi_g(N)
= S_g \otimes_R (S_g^\vee \otimes_{\Gg_g(R)} N)
\cong (S_g \otimes_R S_g^\vee) \otimes_{\Gg_g(R)} N
\cong \operatorname{End}_R(S_g) \otimes_{\Gg_g(R)} N
\cong \Gg_g(R) \otimes_{\Gg_g(R)} N \cong N$
via the coevaluation and the isomorphism
$\rho_g : \Gg_g(R) \xrightarrow{\sim} \operatorname{End}_R(S_g)$.
\end{proof}

\begin{remark}[The inverse functor explicitly]
% label removed: rem:tholog-inverse-functor
Let $N$ be a left $\Gg_g(R)$-module.
As a left $R$-module,
$\Psi_g(N) = S_g^\vee \otimes_{\Gg_g(R)} N$
carries the residual $R$-action from $S_g^\vee$.
In concrete terms: $S_g^\vee$ has basis
$\{e^{\epsilon_1\cdots\epsilon_g}\}$ dual to
$\{e_{\epsilon_1\cdots\epsilon_g}\}$.
The tensor product over $\Gg_g(R)$ imposes the relations
\[
f \cdot \rho_g(c) \otimes n = f \otimes c \cdot n
\qquad
(c \in \Gg_g(R),\; f \in S_g^\vee,\; n \in N),
\]
which identifies $\Psi_g(N)$ with the ``isotypic $e_+$-component''
of $N$ after choosing a decomposition of $N$ as a
$\operatorname{Mat}_{2^g}(R)$-module.
Alternatively, $\Psi_g(N) = e \cdot N$ where
$e = e_{++\cdots+} \cdot e^{++\cdots+} \in \Gg_g(R)$
is the rank-$1$ idempotent $u^{-g} \cdot \prod_{i=1}^g \beta_i\alpha_i$.
\end{remark}

\begin{corollary}[Rank-1 idempotent; \ClaimStatusProvedHere]
% label removed: cor:tholog-rank1-idempotent
In $\Gg_g(R)$, the element
\[
e_g \;:=\; u^{-g} \prod_{i=1}^g \beta_i\alpha_i
\]
(the product taken in any order, since the factors for
distinct handles commute) is an idempotent:
$e_g^2 = e_g$.
The left ideal $\Gg_g(R) \cdot e_g$ is isomorphic to $S_g$
as a left $\Gg_g(R)$-module, and $e_g \cdot \Gg_g(R) \cdot e_g \cong R$.
\end{corollary}

\begin{proof}
Under the isomorphism
$\rho_g : \Gg_g(R) \cong \operatorname{Mat}_{2^g}(R)$,
the element $\beta_i\alpha_i$ maps to
$E_{12}(u\,E_{21}) = u\,E_{11}$
in the $i$-th tensor factor.
(Here we use $E_{12}E_{21} = E_{11}$ in $\operatorname{Mat}_2$.)
Therefore
$\rho_g\!\left(\prod_{i=1}^g \beta_i\alpha_i\right)
= u^g \cdot E_{11}^{\otimes g}$,
and $e_g = u^{-g} \cdot u^g \cdot E_{11}^{\otimes g}
= E_{11}^{\otimes g}$, which is the diagonal matrix with
a single $1$ in the $(1,1)$-position.
This is manifestly an idempotent.
The left ideal $\operatorname{Mat}_{2^g}(R) \cdot E_{11}^{\otimes g}$
is the first column, which is $R^{2^g} = S_g$.
The corner algebra
$E_{11}^{\otimes g} \cdot \operatorname{Mat}_{2^g}(R)
\cdot E_{11}^{\otimes g} = R$ is the $(1,1)$-entry.
\end{proof}

\begin{remark}[Compatibility with genus stabilization]
% label removed: rem:tholog-genus-stabilization
\index{genus stabilization!Morita perspective}
The inclusion $\Gg_g(R) \hookrightarrow \Gg_{g+1}(R)$
(extending by the identity on the $(g{+}1)$-th handle pair)
corresponds under the Morita equivalence to the functor
$M \mapsto M \oplus M$ (doubling of $R$-modules).
This is the algebraic shadow of the genus-stabilization map
$\Sigma_g \hookrightarrow \Sigma_{g+1}$ by
handle attachment: at the module level, each new handle
doubles the rank of the spinor module, reflecting the
doubling of $H_1$ by the new cycle pair.
\end{remark}

\subsubsection{The Morita equivalence as a derived statement}
% label removed: subsubsec:tholog-derived-morita

The Morita equivalence proved above holds at the
abelian/module level, hence \emph{a fortiori} at the derived level.
We record the precise derived statement for completeness.

\begin{theorem}[Derived Morita equivalence; \ClaimStatusProvedHere]
% label removed: thm:tholog-derived-morita
For $R = B_\Theta[u^{-1}]$ and $g \geq 1$, the functors
\[
\Phi_g = S_g \otimes_R^L (-) :
D(R\text{-}\mathbf{mod}) \to D(\Gg_g(R)\text{-}\mathbf{mod})
\]
and
\[
\Psi_g = S_g^\vee \otimes_{\Gg_g(R)}^L (-) :
D(\Gg_g(R)\text{-}\mathbf{mod}) \to D(R\text{-}\mathbf{mod})
\]
are mutually inverse equivalences of triangulated categories.
In particular, for any $\Gg_g(R)$-modules $M, N$:
\[
R\!\operatorname{Hom}_{\Gg_g(R)}(M, N)
\;\simeq\;
R\!\operatorname{Hom}_R(\Psi_g M, \Psi_g N).
\]
\end{theorem}

\begin{proof}
Since $S_g$ is a finitely generated free $R$-module, the
derived tensor product $S_g \otimes_R^L (-)$ coincides with the
underived tensor product.
The unit and counit of the derived adjunction are
quasi-isomorphisms because they are isomorphisms at the
abelian level
(Theorem~\ref{thm:tholog-full-morita-g}(iv)).
\end{proof}

\begin{remark}[The Morita equivalence is canonical]
% label removed: rem:tholog-morita-canonical
The spinor module $S_g$ depends on the choice of symplectic
basis $\{a_1,b_1,\dots,a_g,b_g\}$ for $H_1(\Sigma_g)$.
A change of symplectic basis acts on the Clifford generators
by a symplectic transformation $\gamma \in \operatorname{Sp}_{2g}(\mathbb{Z})$,
which transports the spinor module to an isomorphic one
(because the Clifford algebra $\operatorname{Cl}_{2g}(R,u)$ does
not depend on the basis choice).
The Morita \emph{equivalence class} is therefore independent
of the symplectic basis.
The precise transport is the Weil representation of
$\operatorname{Sp}_{2g}$ on $S_g$: a well-defined projective
representation whose cocycle is controlled by the Maslov index.
\end{remark}

\subsection{The gravitational regime $u = 0$ in full detail}
% label removed: subsec:tholog-gravitational-regime

\index{gravitational regime|textbf}
\index{secondary anomaly!vanishing locus}
\index{exterior algebra!genus completion}
\index{Virasoro algebra!c=26@$c = 26$}

The locus $u = 0$ is where the secondary anomaly
$u = \eta^2$ vanishes, i.e.\ where the transgression generator
becomes nilpotent of order two.
Theorem~\ref{thm:tholog-exterior-degeneration} identified the
quotient $\Gg_g(B_\Theta)/(u)$ as an exterior algebra extension.
The present section analyzes this degenerate regime in full:
the algebra structure, the module theory, and the physical
specialization to Virasoro at $c = 26$.

\subsubsection{The gravitational algebra}
% label removed: subsubsec:tholog-gravitational-algebra

\begin{definition}[Gravitational algebra]
% label removed: def:tholog-gravitational-algebra
The \emph{gravitational algebra} at genus $g$ is
\[
\Gg_g^{\mathrm{grav}}
\;:=\;
\Gg_g(B_\Theta)/(u)
\;\cong\;
B_\Theta/(u) \otimes \Lambda(\alpha_1,\beta_1,\dots,\alpha_g,\beta_g).
\]
\end{definition}

\begin{proposition}[Structure of $B_\Theta/(u)$; \ClaimStatusProvedHere]
% label removed: prop:tholog-bt-mod-u
As a dg algebra,
\[
B_\Theta/(u)
\;\cong\;
\frac{B[\eta]}{(\eta^2, \;\eta b - (-1)^{|b|}b\eta)},
\qquad
|\eta| = 1,\qquad d\eta = \Theta.
\]
The quotient $B_\Theta/(\eta)$ is $B$ itself, and the sequence
\[
0 \to B[-1] \xrightarrow{\cdot\eta} B_\Theta/(u) \to B \to 0
\]
is exact.
If in addition $[\Theta] = 0$ in $H^\bullet(B)$
(i.e.\ if $\Theta = d\xi$ for some $\xi \in B^1$), then
$B_\Theta/(u)$ is quasi-isomorphic to the trivial
square-zero extension $B \oplus B[-1]$.
\end{proposition}

\begin{proof}
The relation $u = \eta^2 = 0$ means $\eta$ squares to zero in
$B_\Theta/(u)$.
Together with the skew-commutation $\eta b = (-1)^{|b|}b\eta$,
this forces $B_\Theta/(u) \cong B \otimes \Lambda(\eta)$
as a graded algebra (but the differential is nontrivial:
$d\eta = \Theta$).

The exact sequence is the short exact sequence of the extension
by $\eta$: the kernel of the projection
$B_\Theta/(u) \to B_\Theta/(\eta) = B$
is the ideal $(\eta) \cong B \cdot \eta \cong B[-1]$
(the shift accounts for $|\eta| = 1$).
Injectivity of multiplication by $\eta$ follows because
$B_\Theta/(u)$ is free of rank $2$ over $B$ with basis
$\{1, \eta\}$.

If $\Theta = d\xi$, define a chain map
$\phi : B_\Theta/(u) \to B \oplus B[-1]$
by $\phi(b_0 + b_1\eta) = (b_0, b_1)$, with the differential
on $B \oplus B[-1]$ being
$d'(b_0, b_1) = (db_0 + \Theta b_1, db_1)
= (db_0 + d\xi \cdot b_1, db_1)$.
Since $\Theta = d\xi$, the map
$\psi(b_0, b_1) = b_0 + b_1(\eta - \xi)$
gives an inverse chain map (checking: $d(\eta - \xi)
= \Theta - d\xi = 0$, and $(\eta - \xi)^2
= \eta^2 - \eta\xi - \xi\eta + \xi^2 = -\eta\xi - \xi\eta + \xi^2$;
since $|\xi| = 1$ and $\eta$ graded-commutes with elements of $B$,
$\eta\xi = (-1)^{|\xi|}\xi\eta = -\xi\eta$, so
$(\eta - \xi)^2 = -\eta\xi + \eta\xi + \xi^2 = \xi^2$,
which vanishes in $B$ if $\xi^2 = 0$, e.g.\ when $B$ is
graded commutative).
In general, $\phi$ is a quasi-isomorphism by the five-lemma
applied to the long exact sequences.

To verify that $\psi$ is indeed inverse to $\phi$: compute
$\phi \circ \psi(b_0, b_1) = \phi(b_0 + b_1(\eta - \xi))
= (b_0, b_1)$ (since the $B$-component of
$b_1(\eta - \xi)$ is $-b_1\xi$ and the $B\eta$-component
is $b_1\eta$), and
$\psi \circ \phi(b_0 + b_1\eta) = \psi(b_0, b_1)
= b_0 + b_1(\eta - \xi)$.
These agree on generators provided $|\xi| = 1$: this holds
because $\xi$ is the image of $\eta$ under the
quasi-isomorphism, and $|\eta| = 1$ by construction.
\end{proof}

\begin{proposition}[Exterior algebra basis; \ClaimStatusProvedHere]
% label removed: prop:tholog-exterior-basis
The gravitational algebra $\Gg_g^{\mathrm{grav}}$
has a $B$-module basis consisting of all monomials
\begin{equation}% label removed: eq:exterior-monomial-basis
\eta^{\epsilon_0}\,
\alpha_1^{a_1}\beta_1^{b_1}
\cdots
\alpha_g^{a_g}\beta_g^{b_g},
\qquad
\epsilon_0 \in \{0,1\},\;
a_i, b_i \in \{0,1\}.
\end{equation}
There are $2 \cdot 4^g = 2^{2g+1}$ such monomials.
The $B$-module rank is $2^{2g+1}$.
\end{proposition}

\begin{proof}
In $\Gg_g^{\mathrm{grav}}$, the generators $\eta, \alpha_i, \beta_i$
are all odd (degree $1$) and all square to zero:
$\eta^2 = u = 0$, $\alpha_i^2 = 0$, $\beta_i^2 = 0$.
All pairs anticommute: this includes the
cross-handle relations (which are zero in $\Gg_g$ already)
and the Clifford relations
$\alpha_i\beta_j + \beta_j\alpha_i = \delta_{ij}u = 0$.
Also, $\eta$ anticommutes with each $\alpha_i$ and $\beta_i$
because $\eta$ graded-commutes with elements of $B_\Theta$
and $|\alpha_i| = 1$ is odd:
$\eta\alpha_i = (-1)^{|\alpha_i|}\alpha_i\eta = -\alpha_i\eta$.

Therefore all $2g + 1$ generators $\{\eta, \alpha_1, \beta_1,
\dots, \alpha_g, \beta_g\}$ mutually anticommute and square to
zero.
This is exactly the exterior algebra on $2g+1$ generators.
The monomials~\eqref{eq:exterior-monomial-basis} span, and
the argument of Lemma~\ref{lem:tholog-clifford-pbw} (reducing
modulo $u$ and $\Theta$) shows they are linearly independent
over $B$.
\end{proof}

\begin{corollary}[Gravitational algebra as exterior algebra; \ClaimStatusProvedHere]
% label removed: cor:tholog-gravitational-exterior
\[
\Gg_g^{\mathrm{grav}}
\;\cong\;
B_\Theta/(u) \otimes \Lambda^{2g}(\alpha_1,\beta_1,\dots,\alpha_g,\beta_g)
\;\cong\;
B \otimes \Lambda^{2g+1}(\eta,\alpha_1,\beta_1,\dots,\alpha_g,\beta_g)
\]
as graded algebras, where $\Lambda^{2g+1}$ denotes the exterior
algebra on $2g + 1$ generators of degree~$1$.
The differential on the right-hand side acts by
$d\eta = \Theta$, $d\alpha_i = d\beta_i = 0$, and extends
to monomials by the Leibniz rule.
\end{corollary}

\begin{remark}[Topological interpretation]
% label removed: rem:tholog-grav-topological
The $2g + 1$ exterior generators decompose as:
\begin{itemize}
\item $\alpha_1, \beta_1, \dots, \alpha_g, \beta_g$: these
  are the parity-shifted homology classes of
  $H_1(\Sigma_g; k)$, generating $\Lambda^\bullet(H_1(\Sigma_g))$.
\item $\eta$: this is the transgression generator, representing
  the fundamental class of the circle
  $S^1 \subset \Sigma_g$ (the ``time circle'' in the
  physics interpretation).
\end{itemize}
The exterior algebra $\Lambda^{2g+1}$ is therefore
$\Lambda^\bullet\bigl(H_1(\Sigma_g) \oplus k\langle\eta\rangle\bigr)$.
When $\Theta = 0$, the differential is trivial and the
cohomology of $\Gg_g^{\mathrm{grav}}$ is
$H^\bullet(B) \otimes \Lambda^{2g+1}$.
\end{remark}

\begin{proposition}[Module theory on the gravitational locus; \ClaimStatusProvedHere]
% label removed: prop:tholog-grav-modules
Let $M$ be a left $\Gg_g^{\mathrm{grav}}$-module.
Then $M$ decomposes (not canonically) as
\[
M \;\cong\; M_0 \otimes \Lambda^{2g+1}
\]
where $M_0 := M / (\eta, \alpha_1, \beta_1, \dots, \alpha_g, \beta_g) \cdot M$
is the ``top component'' and $\Lambda^{2g+1}$ acts freely.
This decomposition holds if and only if all $2g + 1$ generators
act injectively on the quotient $M / (\text{previous generators})$
at each stage.

More precisely: $M$ is a module over $B \otimes \Lambda^{2g+1}$
with differential $d_M + \Theta \cdot \iota_\eta$,
where $\iota_\eta$ is contraction by $\eta$ in the exterior factor.
\end{proposition}

\begin{proof}
Since $\Gg_g^{\mathrm{grav}} = B \otimes \Lambda^{2g+1}$
(Corollary~\ref{cor:tholog-gravitational-exterior}),
a $\Gg_g^{\mathrm{grav}}$-module is a
$B \otimes \Lambda^{2g+1}$-module.
The $\Lambda^{2g+1}$-module structure is a collection of
$2g + 1$ mutually anticommuting, square-zero endomorphisms
acting $B$-linearly.
The differential incorporates the term $d\eta = \Theta$
via $\iota_\eta$ contraction.
The decomposition $M \cong M_0 \otimes \Lambda^{2g+1}$
holds whenever $M$ is free over $\Lambda^{2g+1}$,
which is the generic case (the exterior algebra over a field
of characteristic $0$ is self-injective, so free = projective
= injective for finitely generated modules).
\end{proof}

\subsubsection{Virasoro at $c = 26$: effective scalar cancellation}
% label removed: subsubsec:tholog-virasoro-c26

The Virasoro algebra at $c = 26$ is the unique member of the
Virasoro family where the gravitational regime is physically
realized.

\begin{proposition}[Virasoro $c = 26$: vanishing effective scalar curvature; \ClaimStatusProvedHere]
% label removed: prop:tholog-virasoro-c26
\index{Virasoro algebra!c=26@$c = 26$!gravitational regime}
For $\mathrm{Vir}_{26}$:
\begin{enumerate}[label=\textup{(\roman*)}]
\item The intrinsic modular Koszul curvature is
  $\kappa(\mathrm{Vir}_{26}) = 26/2 = 13$.
  The \emph{effective} curvature, after ghost subtraction, vanishes:
  $\kappa_{\mathrm{eff}} = \kappa(\mathrm{Vir}_{26}) - 13
  = (c-26)/2\big|_{c=26} = 0$.
\item The Koszul dual is
  $\mathrm{Vir}_{26}^! \simeq \mathrm{Vir}_0$
  (the Virasoro at $c = 0$, whose scalar curvature vanishes
  but which is not the zero algebra).
\item The effective scalar anomaly class vanishes:
  $\Theta^{\min}_{\mathrm{eff}} = \kappa_{\mathrm{eff}} \cdot \omega_1 = 0$.
\item The effective secondary scalar anomaly vanishes:
  $u_{\mathrm{eff}} = \kappa_{\mathrm{eff}} \cdot \omega_g = 0$.
\item These scalar cancellations do not by themselves identify the
  full matter--ghost transgression algebra or the full coupled genus
  completion.
\end{enumerate}
\end{proposition}

\begin{proof}
Part~(i): by Vol~I, Theorem~D, $\kappa(\mathrm{Vir}_c) = c/2$
for all~$c$.  At $c = 26$: $\kappa = 13$.  The effective
curvature is $\kappa_{\mathrm{eff}} = (c-26)/2 = 0$,
reflecting matter-ghost cancellation (not vanishing of the
intrinsic curvature).

Part~(ii) is the Virasoro duality relation
$\mathrm{Vir}_{26}^! \simeq \mathrm{Vir}_0$ together with the
already-established fact that $c=0$ is the uncurved scalar locus, not
the zero algebra.  Parts~(iii)--(iv) are the corresponding scalar
vanishing statements.  Part~(v) is the scope warning forced by the
non-scalar Virasoro shadow obstruction tower.
\end{proof}

\begin{theorem}[Zero-anomaly specialization of the transgression algebra; \ClaimStatusProvedHere]
% label removed: thm:tholog-gravitational-collapse-c26
For any formal specialization with $\Theta = 0$,
the genus-$g$ Clifford completion is
\[
\Gg_g\bigl(B_{\Theta=0}\bigr)
\;\cong\;
B \otimes k[\eta] \otimes \Lambda^{2g}(\alpha_1,\beta_1,\dots,\alpha_g,\beta_g)
/\bigl(\alpha_i\beta_j + \beta_j\alpha_i = \delta_{ij}\eta^2\bigr).
\]
Setting $u = \eta^2$:
\begin{enumerate}[label=\textup{(\roman*)}]
\item The Clifford algebra is \emph{not} automatically exterior
  (because $u = \eta^2 \neq 0$ as an element of $k[\eta]$).
\item Modulo $u$: the algebra is
  $B \otimes \Lambda^{2g+1}(\eta, \alpha_1, \beta_1, \dots,
  \alpha_g, \beta_g)$
  with trivial differential ($d\eta = \Theta = 0$).
  This is the full exterior algebra on $H_1(\Sigma_g) \oplus k\langle\eta\rangle$.
\item After inverting $u$: the algebra is Morita-trivial,
  $\Gg_g[u^{-1}] \cong \operatorname{Mat}_{2^g}(B[\eta, \eta^{-1}])$.
\end{enumerate}
This is a formal zero-anomaly specialization statement.  For the
critical bosonic string at $c = 26$, the proved input is only the
effective scalar cancellation $\kappa_{\mathrm{eff}} = 0$; identifying
the full coupled transgression algebra with this specialization is not
proved here.
\end{theorem}

\begin{proof}
Since $\Theta = 0$, the transgression algebra is
$B_\Theta = B \otimes k[\eta]$ with $d\eta = 0$.
The Clifford relations remain nontrivial:
$\alpha_i\beta_j + \beta_j\alpha_i = \delta_{ij}\eta^2$.
Since $\eta^2$ is a nonzero element of $k[\eta]$
(the polynomial ring), the Clifford algebra over
$B[\eta]$ does not automatically degenerate.
Setting $u = \eta^2 = 0$ forces the exterior degeneration
(Theorem~\ref{thm:tholog-exterior-degeneration}).
Inverting $u$ gives the Morita triviality
(Theorem~\ref{thm:tholog-morita-triviality}).
The differential on the exterior algebra modulo $u$ is
$d\eta = 0$, $d\alpha_i = d\beta_i = 0$, and $dB = dB$:
the only nontrivial differential is that of $B$ itself.
\end{proof}

\begin{remark}[Comparison with BRST cohomology at $c = 26$]
% label removed: rem:tholog-brst-c26
\index{BRST cohomology!c=26@$c = 26$}
The physical significance of $c = 26$ is that the ghost
BRST system $(b, c)$ at ghost number anomaly
$c_{\mathrm{gh}} = -26$ cancels the matter Virasoro anomaly
$c = 26$, giving total central charge $0$.
In the anomaly-completed framework:
\begin{itemize}
\item The ghost system contributes the exterior factor
  $\Lambda(b, c)$ with $|b| = 2$, $|c| = -1$, and the
  BRST differential $Q_{\mathrm{BRST}}$ acts nontrivially
  on this factor.
\item The matter--ghost coupled system at $c_{\mathrm{total}} = 0$
  has vanishing effective scalar characteristic
  $\kappa_{\mathrm{eff}} = 0$.  The vanishing of the full effective
  anomaly $\Theta_{\mathrm{eff}} = 0$ requires the complete
  matter--ghost cancellation at all genera and all orders of the
  shadow obstruction tower, not merely the scalar projection
  $\kappa_{\mathrm{eff}} = 0$.  Identifying the full coupled genus
  completion with the zero-anomaly specialization above is not proved
  here.
\item The BRST cohomology of the genus-$g$ exterior algebra
  $\Lambda^{2g+1}$ supplies only the ghost zero-mode bookkeeping
  suggested by the algebraic shadow.  Identifying it with the
  full genus-$g$ physical state space, or with the $g$-th
  coefficient of the bosonic-string partition function, requires
  the full coupled BRST/bar comparison and is not proved here.
\item The zero-anomaly specialization above is the algebraic model
  one would obtain if the full coupled transgression class vanished.
  The present manuscript proves only the scalar cancellation needed
  for the genus-$1$ anomaly budget.
\end{itemize}
\end{remark}

\subsubsection{Cohomology of the gravitational algebra}
% label removed: subsubsec:tholog-grav-cohomology

\begin{proposition}[Cohomology when $\Theta \neq 0$; \ClaimStatusProvedHere]
% label removed: prop:tholog-grav-cohomology
Suppose $[\Theta] \neq 0$ in $H^\bullet(B)$
(the generic case for $c \neq 26$ in the Virasoro family).
Then:
\[
H^\bullet(\Gg_g^{\mathrm{grav}})
\;\cong\;
\frac{H^\bullet(B)}{([\Theta])}
\;\otimes\;
\Lambda^{2g}(\alpha_1,\beta_1,\dots,\alpha_g,\beta_g).
\]
In particular, the cohomology is a quotient of the
exterior algebra by the principal ideal generated by the
class $[\Theta]$.
\end{proposition}

\begin{proof}
By Corollary~\ref{cor:tholog-gravitational-exterior},
\[
\Gg_g^{\mathrm{grav}}
= B \otimes \Lambda(\eta) \otimes \Lambda(\alpha_1,\dots,\beta_g).
\]
The differential acts only on the $B \otimes \Lambda(\eta)$
factor (by $d\eta = \Theta$, $d\alpha_i = d\beta_i = 0$).
By the K\"{u}nneth theorem,
\[
H^\bullet(\Gg_g^{\mathrm{grav}})
= H^\bullet(B \otimes \Lambda(\eta)) \otimes \Lambda^{2g}.
\]
The complex $B \otimes \Lambda(\eta)$ with $d\eta = \Theta$
is the two-term complex $B \xrightarrow{\cdot\Theta} B[1]$
(the short exact sequence from
Proposition~\ref{prop:tholog-bt-mod-u}).
Its cohomology is
$\ker(\cdot\Theta) \oplus \operatorname{coker}(\cdot\Theta)[-1]$.
If $[\Theta] \neq 0$ and acts injectively on
$H^\bullet(B)$, then the long exact sequence from
Theorem~\ref{thm:tholog-hypersurface} specialized to
$u = 0$ gives $H^\bullet(B \otimes \Lambda(\eta))
\cong H^\bullet(B)/([\Theta])$.
\end{proof}

\begin{proposition}[Cohomology when $\Theta = 0$; \ClaimStatusProvedHere]
% label removed: prop:tholog-grav-cohomology-zero
When $\Theta = 0$:
\[
H^\bullet(\Gg_g^{\mathrm{grav}})
\;\cong\;
H^\bullet(B) \otimes \Lambda^{2g+1}(\eta,\alpha_1,\beta_1,\dots,\alpha_g,\beta_g).
\]
\end{proposition}

\begin{proof}
With $\Theta = 0$, the differential on $\Lambda(\eta)$ is zero,
and K\"{u}nneth gives the result.
\end{proof}

\subsubsection{The Hodge decomposition on the gravitational locus}
% label removed: subsubsec:tholog-hodge-grav

\begin{proposition}[Hodge-type decomposition; \ClaimStatusProvedHere]
% label removed: prop:tholog-hodge-grav
The exterior algebra
$\Lambda^{2g+1}(\eta, \alpha_1, \beta_1, \dots, \alpha_g, \beta_g)$
carries a bigrading:
\begin{itemize}
\item $p = \deg_\eta \in \{0, 1\}$ (``transgression degree'');
\item $q = \deg_{\alpha,\beta} \in \{0, 1, \dots, 2g\}$
  (``handle degree'').
\end{itemize}
The total exterior degree is $p + q$.
The gravitational algebra decomposes as
\[
\Gg_g^{\mathrm{grav}}
= \bigoplus_{p=0}^{1} \bigoplus_{q=0}^{2g}
\Gg_g^{p,q},
\qquad
\Gg_g^{p,q} = B^{p,q} \otimes \Lambda^p(\eta) \otimes
\Lambda^q(\alpha_1,\beta_1,\dots,\alpha_g,\beta_g),
\]
and the differential maps $\Gg_g^{0,q} \to \Gg_g^{1,q}$
(by $d\eta = \Theta$) and $\Gg_g^{1,q} \to \Gg_g^{0,q}$
trivially (since $\Theta \cdot \eta = 0$ in
$\Gg_g^{\mathrm{grav}}$ because $\eta^2 = 0$ and
$\Theta = d\eta$).
\end{proposition}

\begin{proof}
This is a direct decomposition by monomial type.
The differential $d$ on $\eta$ sends $1 \mapsto 0$ and
$\eta \mapsto \Theta$ (which lies in $B$, hence in $p = 0$).
The generators $\alpha_i, \beta_i$ are $d$-closed, so the
differential does not change $q$.

It remains to verify that $\Theta \cdot \eta = 0$ in
$\Gg_g^{\mathrm{grav}}$, so that the differential does not
map from $p = 1$ back to $p = 0$ nontrivially.
In the quotient $\Gg_g^{\mathrm{grav}}$ by $u = \eta^2$,
compute $d(u) = d(\eta^2)$.  By the Leibniz rule,
$d(\eta^2) = (d\eta) \cdot \eta + \eta \cdot (d\eta)
= \Theta \cdot \eta + \eta \cdot \Theta$.
By the skew-commutation rule
$\eta \cdot b = (-1)^{|b|} b \cdot \eta$ for $b \in B$,
and since $|\Theta| = 2$ is even, we get
$\eta \cdot \Theta = (-1)^{2} \Theta \cdot \eta
= \Theta \cdot \eta$.
Hence $d(u) = 2\,\Theta \cdot \eta$.
But $u = 0$ in $\Gg_g^{\mathrm{grav}}$, so
$d(u) = 0$, giving $\Theta \cdot \eta = 0$
(in characteristic $\neq 2$).
\end{proof}

\subsection{Critical points for the standard landscape}
% label removed: subsec:tholog-critical-points

\index{critical point!anomaly completion|textbf}
\index{self-dual point}
\index{secondary anomaly!standard families}

The critical and self-dual points of the standard families are
the loci where the anomaly-completed structure simplifies:
at self-dual points the curvature vanishes, and the genus tower
collapses to the gravitational regime; at critical levels
the curvature diverges or degenerates.
The following subsection computes these special points for each
standard family and records the secondary anomaly hierarchy.

\subsubsection{Virasoro}
% label removed: subsubsec:tholog-critical-virasoro

\begin{computation}[Virasoro critical loci; \ClaimStatusProvedHere]
% label removed: comp:tholog-virasoro-critical
\index{Virasoro algebra!critical points}
The Virasoro family $\mathrm{Vir}_c$ has:
\begin{enumerate}[label=\textup{(\roman*)}]
\item \emph{Critical dimension}: $c = 26$, where
  $\kappa(\mathrm{Vir}_{26}) = 13$ but the effective curvature
  $\kappa_{\mathrm{eff}} = (c-26)/2 = 0$
  (see the effective scalar cancellation discussion above).
  The effective scalar anomaly class vanishes; this does not by itself
  identify the full coupled genus tower with the zero-anomaly
  specialization.
\item \emph{Self-dual point}: $c = 13$, where
  $\mathrm{Vir}_{13}^! \simeq \mathrm{Vir}_{13}$
  (Vol~I, Theorem~A).
  Here $\kappa(\mathrm{Vir}_{13}) = 13/2$; the complementarity
  sum is $\kappa(\mathrm{Vir}_{13}) + \kappa(\mathrm{Vir}_{13})
  = 13$, reflecting the non-vanishing constant
  in the Virasoro complementarity formula.
\item \emph{Distinction between critical and self-dual}:
  $c = 26$ is the locus where
  $\kappa_{\mathrm{eff}} = 0$ (matter-ghost cancellation);
  $c = 13$ is the locus of self-duality
  ($\kappa(\mathrm{Vir}_{13}) = 13/2$, nonzero but
  symmetric under duality).
  The effective anomaly vanishes at $c = 26$; at $c = 13$
  it takes the value $\kappa_{\mathrm{eff}} = -13/2$.
\item \emph{Quartic resonance}: the contact quartic invariant is
  $Q_{\mathrm{Vir}}^{\mathrm{contact}}
  = \frac{10}{c(5c + 22)}$
  (Vol~I; see also
  Computation~\ref{comp:virasoro-anomaly-completion} above).
  This has poles at $c = 0$ (the uncurved scalar Virasoro point) and
  $c = -22/5$ (a minimal model value).
  Its zeros: $Q = 0$ iff $c = \infty$ (decompactification limit).
\end{enumerate}
\end{computation}

\subsubsection{$\mathcal{W}_3$}
% label removed: subsubsec:tholog-critical-w3

\begin{computation}[$\mathcal{W}_3$ critical loci; \ClaimStatusProvedHere]
% label removed: comp:tholog-w3-critical
\index{W-algebra@$\mathcal W$-algebra!W3@$\mathcal W_3$!critical points}
The $\mathcal{W}_3$ algebra has generators $T$ (weight $2$)
and $W$ (weight $3$), with Koszul duality
$\mathcal{W}_{3,c}^! \simeq \mathcal{W}_{3,100-c}$
(Vol~I, Theorem~A; the critical sum is $c + c' = 100$).

\begin{enumerate}[label=\textup{(\roman*)}]
\item \emph{Critical dimension}: $c = 100$, where the Koszul dual
  $\mathcal{W}_{3,0}$ is trivial with
  $\kappa(\mathcal{W}_{3,0}) = 0$.
  The algebra itself has $\kappa(\mathcal{W}_{3,100}) = 250/3$;
  the anomaly-completed datum degenerates because the
  \emph{dual} genus tower collapses (as for
  $\mathrm{Vir}_{26}$, where $\kappa(\mathrm{Vir}_0) = 0$).
\item \emph{Self-dual point}: $c = 50$, where
  $\mathcal{W}_{3,50}^! \simeq \mathcal{W}_{3,50}$
  and $\kappa(\mathcal{W}_{3,50}) = 125/3$
  (nonzero: the complementarity sum is
  $\kappa + \kappa^! = 250/3$, not~$0$).
\item \emph{The shadow obstruction tower}: type $M$ (infinite depth,
  like Virasoro).
  The cubic shadow is generically nonvanishing; the
  quartic resonance class is a rational function of $c$.
\item \emph{Physical realization}: The $\mathcal{W}_3$
  string at $c = 100$ is the higher-spin analogue of the
  bosonic string at $c = 26$.  The BRST system requires a
  spin-$3$ ghost pair $(b_3, c_3)$ with
  $c_{\mathrm{gh},3} = -74$ in addition to the standard
  $(b, c)$ with $c_{\mathrm{gh}} = -26$.
  Total ghost central charge: $-26 + (-74) = -100$,
  cancelling the matter $c = 100$.
\end{enumerate}
\end{computation}

\subsubsection{$\mathcal{W}_N$ general formula}
% label removed: subsubsec:tholog-critical-wn

\begin{computation}[$\mathcal{W}_N$ critical loci; \ClaimStatusProvedHere]
% label removed: comp:tholog-wn-critical
\index{W-algebra@$\mathcal W$-algebra!WN@$\mathcal W_N$!critical points}
The $\mathcal{W}_N$ algebra is the principal DS reduction of
$\widehat{\mathfrak{sl}}_N$.
It has generators of conformal weights $2, 3, \dots, N$, with
central-charge formula
\[
c(\mathcal{W}_N, k)
= (N-1) - N(N^2-1)\frac{(k+N-1)^2}{k+N}.
\]
The Koszul duality is
$\mathcal{W}_{N,c}^! \simeq \mathcal{W}_{N,c_N^*-c}$
where the critical sum is
\begin{equation}% label removed: eq:wn-critical-sum
c_N^* \;=\; 2(N-1)(2N^2 + 2N + 1).
\end{equation}
This formula is computed from the BRST ghost system:
for principal reduction, the ghost system has
$\dim(\mathfrak{m}^+) = N(N-1)/2$ fermionic pairs
$(b_j, c_j)$ for $j = 2, \dots, N$, with total ghost
central charge
\[
c_{\mathrm{gh}}^{\mathcal{W}_N}
= -2\sum_{j=2}^{N}(6j^2 - 6j + 1)
= -2(N-1)(2N^2 + 2N + 1).
\]
The critical sum $c_N^*$ equals $|c_{\mathrm{gh}}|$.

\begin{center}
\renewcommand{\arraystretch}{1.3}
\begin{tabular}{cccc}
$N$ & $c_N^*$ (critical sum) & $c_N^*/2$ (self-dual) &
$c_N^*$ (critical dim.) \\
\hline
$2$ & $26$ & $13$ & $26$ \\
$3$ & $100$ & $50$ & $100$ \\
$4$ & $246$ & $123$ & $246$ \\
$5$ & $488$ & $244$ & $488$ \\
$6$ & $850$ & $425$ & $850$ \\
$N$ & $2(N{-}1)(2N^2{+}2N{+}1)$ & $(N{-}1)(2N^2{+}2N{+}1)$ &
$2(N{-}1)(2N^2{+}2N{+}1)$
\end{tabular}
\end{center}

\emph{Properties.}
\begin{enumerate}[label=\textup{(\roman*)}]
\item The critical dimension $c_N^*$ grows as $2N^3$ for
  large $N$.
\item The self-dual point $c_N^*/2$ is half-integral for
  even $N$ (because $c_N^*$ is odd when $N$ is even).
\item At $c = c_N^*$: the Koszul dual $\mathcal{W}_{N,0}$
  is trivial with $\kappa(\mathcal{W}_{N,0}) = 0$, and
  the dual genus tower collapses.  The algebra itself
  has $\kappa(\mathcal{W}_{N,c_N^*}) = (H_N{-}1)\,c_N^* \neq 0$.
\item At $c = c_N^*/2$: the algebra is self-dual,
  $\kappa = (H_N{-}1)\,c_N^*/2$ (nonzero for $N \ge 2$;
  the complementarity sum $\kappa + \kappa^! = (H_N{-}1)\,c_N^*$
  is nonzero, unlike the KM case).
\item The quartic resonance class
  $Q^{\mathrm{contact}}_{\mathcal{W}_N}$ is a rational
  function of $c$ whose poles are at $c = 0$ and at
  finitely many minimal-model values.
\end{enumerate}
\end{computation}

\subsubsection{Affine Kac--Moody}
% label removed: subsubsec:tholog-critical-affine

\begin{computation}[Affine critical loci; \ClaimStatusProvedHere]
% label removed: comp:tholog-affine-critical
\index{Kac--Moody algebra!affine!critical points}
\index{critical level}
The affine Kac--Moody algebra $\widehat{\mathfrak{g}}_k$
at level $k$ has central charge
\[
c(\widehat{\mathfrak{g}}_k)
= \frac{k \dim\mathfrak{g}}{k + h^\vee},
\]
and chiral Koszul duality gives
$\kappa(\widehat{\mathfrak{g}}_k^!)
= \kappa(\widehat{\mathfrak{g}}_{-k-2h^\vee})$
(the Koszul dual is the chiral CE algebra at the Feigin--Frenkel dual level).

\begin{enumerate}[label=\textup{(\roman*)}]
\item \emph{Critical level}: $k = -h^\vee$, where the central
  charge $c(\widehat{\mathfrak{g}}_{-h^\vee})$ is undefined
  (the Sugawara construction fails: the Segal--Sugawara operator
  $S = \frac{1}{2(k+h^\vee)}\sum {:}J^a J^a{:}$ has a pole).
  The modular Koszul curvature
  $\kappa(\widehat{\mathfrak{g}}_{-h^\vee})
  = 0$ (the bar complex is flat), but the anomaly-completed
  datum is ill-defined because it requires the Sugawara tensor.
  The critical level is a genuine singularity of the
  Sugawara construction, not merely a zero of~$\kappa$.

\item \emph{Self-dual level}: $k + h^\vee = -(k + h^\vee)$,
  i.e.\ $k = -h^\vee$, which is the critical level.
  For affine algebras, the self-dual and critical levels
  \emph{coincide}.
  This is a fundamental difference from the Virasoro/W-algebra
  families, where the self-dual and critical points are distinct.

\item \emph{Central-charge sum}: for generic $k$,
  \[
  c(\widehat{\mathfrak{g}}_k) + c(\widehat{\mathfrak{g}}_{-k-2h^\vee})
  = \frac{k\dim\mathfrak{g}}{k+h^\vee}
  + \frac{(-k-2h^\vee)\dim\mathfrak{g}}{-k-h^\vee}
  = \frac{k\dim\mathfrak{g}}{k+h^\vee}
  + \frac{(k+2h^\vee)\dim\mathfrak{g}}{k+h^\vee}
  = 2\dim\mathfrak{g}.
  \]
  The critical sum is $c_{\mathrm{sum}} = 2\dim\mathfrak{g}$,
  achieved at $c = \dim\mathfrak{g}$ (the ``free field''
  value, which corresponds to $k = \infty$).

\item \emph{Curvature formula}: the curvature
  $\kappa(\widehat{\mathfrak{g}}_k)$ is computed from the
  Killing form on $\mathfrak{g}$ and the level:
  \[
  \kappa(\widehat{\mathfrak{g}}_k)
  = \frac{(k + h^\vee)\dim\mathfrak{g}}{2h^\vee}.
  \]
  This is anti-symmetric: $\kappa(k) = -\kappa(-k-2h^\vee)$
  (the Feigin--Frenkel symmetry).
  It vanishes at $k = -h^\vee$ (critical level) and grows
  linearly in $k + h^\vee$.
\end{enumerate}
\end{computation}

\begin{computation}[$\widehat{\mathfrak{sl}}_2$ explicit; \ClaimStatusProvedHere]
% label removed: comp:tholog-sl2-affine-explicit
\index{Kac--Moody algebra!sl2@$\widehat{\mathfrak{sl}}_2$!critical data}
For $\widehat{\mathfrak{sl}}_2$ with $h^\vee = 2$,
$\dim\mathfrak{g} = 3$:
\begin{center}
\renewcommand{\arraystretch}{1.2}
\begin{tabular}{ll}
Central charge & $c = 3k/(k+2)$ \\
Critical level & $k = -2$ \\
Central-charge sum & $c_{\mathrm{sum}} = 6$ \\
Self-dual central charge & undefined at critical level $k = -h^\vee = -2$ (Sugawara undefined) \\
Curvature & $\kappa = 3(k+2)/4$ \\
Koszul dual & $\widehat{\mathfrak{sl}}_2{}_{-k-4}$ \\
Secondary anomaly & $u = \eta^2$, with $d\eta = \kappa \cdot \omega_1$ \\
Shadow depth & $3$ (type $L$: tree-level, terminates at cubic)
\end{tabular}
\end{center}
\end{computation}

\subsubsection{The master table}
% label removed: subsubsec:tholog-master-table

\begin{computation}[Critical points for all standard families; \ClaimStatusProvedHere]
% label removed: comp:tholog-master-critical-table
\index{critical point!master table}

\begin{center}
\renewcommand{\arraystretch}{1.3}
\begin{tabular}{lccccl}
\textbf{Family} & $c_{\mathrm{sum}}$ & $c_{\mathrm{sd}}$ &
$c_{\mathrm{crit}}$ & $\kappa = 0?$ & \textbf{Shadow depth} \\
\hline
$\cH_k$ (Heisenberg) & $0$ & $0$ & $0$ & always $\kappa = k$
& $2$ (type $G$) \\
$\widehat{\mathfrak{sl}}_2{}_k$ & $6$ & $3$ & $k = -2$ &
at critical & $3$ (type $L$) \\
$\widehat{\mathfrak{sl}}_N{}_k$ & $2\dim\mathfrak{sl}_N$
& $\dim\mathfrak{sl}_N$ & $k = -N$ & at critical &
$3$ (type $L$) \\
$\widehat{\mathfrak{g}}_k$ (general) & $2\dim\mathfrak{g}$
& $\dim\mathfrak{g}$ & $k = -h^\vee$ & at critical &
$3$ (type $L$) \\
$\beta\gamma$ & $-2$ & $-1$ & --- & never ($\kappa \neq 0$)
& $4$ (type $C$) \\
$\mathrm{Vir}_c$ & $26$ & $13$ & $c = 26$ & $\kappa^! = 0$ at $c = 26$
& $\infty$ (type $M$) \\
$\mathcal{W}_{3,c}$ & $100$ & $50$ & $c = 100$ &
$\kappa^! = 0$ at $c = 100$ & $\infty$ (type $M$) \\
$\mathcal{W}_{N,c}$ & $c_N^*$ & $c_N^*/2$ & $c = c_N^*$ &
$\kappa^! = 0$ at $c = c_N^*$ & $\infty$ (type $M$) \\
$\mathcal{B}^k$ (BP) & $196$ & $98$ & --- &
$\kappa^! = 0$ at $c = 196$ & $\infty$ (type $M$)
\end{tabular}
\end{center}

Here $c_N^* = 2(N-1)(2N^2+2N+1)$.
``$\kappa = 0?$'' records where the curvature vanishes.
For KM, $\kappa = 0$ at the critical level (self-dual = critical).
For Virasoro/$\mathcal{W}$-algebras, the \emph{algebra's own}
$\kappa \neq 0$ at the self-dual point; only the
\emph{dual's} curvature $\kappa^!$ vanishes at $c = c_{\mathrm{sum}}$.
\end{computation}

\subsubsection{The secondary anomaly hierarchy}
% label removed: subsubsec:tholog-secondary-hierarchy

\begin{definition}[Secondary anomaly hierarchy]
% label removed: def:tholog-secondary-hierarchy
\index{secondary anomaly!hierarchy}
Given a chiral algebra $\cA$ with anomaly-completed datum
$(B, \Theta, B_\Theta)$, define the \emph{secondary anomaly
hierarchy} as the sequence of central classes in $B_\Theta$:
\begin{align*}
u_1 &:= \eta^2 \in Z^2(B_\Theta) \cap Z(B_\Theta), \\
u_2 &:= \eta^2 \cdot \Theta \in Z^4(B_\Theta) \cap Z(B_\Theta), \\
u_r &:= \eta^2 \cdot \Theta^{r-1}
\in Z^{2r}(B_\Theta) \cap Z(B_\Theta).
\end{align*}
Each $u_r$ is closed (because $u_1$ and $\Theta$ are both
closed and central) and central (because $u_1$ and $\Theta$
are both central).
\end{definition}

\begin{proposition}[Hierarchy filtration; \ClaimStatusProvedHere]
% label removed: prop:tholog-hierarchy-filtration
Assume $[\Theta]^m = 0$ in $H^\bullet(B)$ for some $m \geq 1$
\textup{(}i.e.\ the cohomology class of~$\Theta$ is nilpotent
of order~$m$\textup{)}.
The secondary anomaly hierarchy satisfies:
\begin{enumerate}[label=\textup{(\roman*)}]
\item $u_r = u_1 \cdot \Theta^{r-1}$ for all $r \geq 1$.
\item Under the nilpotence hypothesis $[\Theta]^m = 0$,
  $[u_r] = 0$ in $H^\bullet(B_\Theta)$ for $r \geq m + 1$.
\item In $H^\bullet(B_\Theta)/([\Theta])$, all $u_r$ with
  $r \geq 2$ vanish: $[u_r] = [u_1] \cdot [\Theta]^{r-1} = 0$.
\item The graded ideal $(u_1, u_2, \dots)$ is the principal
  ideal $(u_1)$, since $u_r = u_1 \cdot \Theta^{r-1}$.
\end{enumerate}
\end{proposition}

\begin{proof}
Part~(i) is the definition.
Part~(ii): if $[\Theta]^m = 0$, then $[u_r] = [u_1] \cdot [\Theta]^{r-1}$
vanishes for $r - 1 \geq m$.
Part~(iii): in the quotient by $([\Theta])$, each
$[\Theta]^{r-1} = 0$ for $r \geq 2$.
Part~(iv): $u_r = u_1 \cdot \Theta^{r-1} \in (u_1)$.
\end{proof}

\begin{remark}[The secondary hierarchy and the shadow obstruction tower]
% label removed: rem:tholog-hierarchy-shadow
\index{shadow tower!secondary hierarchy}
The secondary anomaly hierarchy $u_1, u_2, \dots$ is the
$B_\Theta$-internal shadow of the genus-$g$ shadow obstruction
tower from Vol~I.
The arity-$r$ shadow $\mathrm{Sh}_r(\Theta_\cA)$ of the
full MC element projects to $u_r$ in the transgression algebra:
\begin{center}
\renewcommand{\arraystretch}{1.2}
\begin{tabular}{lll}
\textbf{Shadow level} & \textbf{Vol~I object} &
\textbf{Transgression shadow} \\
\hline
$r = 2$ (curvature) & $\kappa(\cA)$ & $u_1 = \eta^2$ \\
$r = 3$ (cubic) & $C(\cA)$ & $u_2 = \eta^2 \Theta$ \\
$r = 4$ (quartic) & $Q(\cA)$ & $u_3 = \eta^2 \Theta^2$
\end{tabular}
\end{center}
The shadow depth classification (G/L/C/M from Vol~I) is
visible in the hierarchy: for type $G$ (Heisenberg),
$\Theta = k/z$ and $\Theta^r$ scales homogeneously, so the
hierarchy does not terminate but is uniformly controlled by
the single parameter $k$.
For type $M$ (Virasoro), each $u_r$ contributes independently
to the module theory.
\end{remark}

\subsection{The symplectic polarization of complementarity}
% label removed: subsec:tholog-symplectic-polarization

\index{complementarity!symplectic polarization|textbf}
\index{Verdier involution}
\index{Lagrangian!complementarity}
\index{cross-pairing}

Volume~I's Theorem~C (the complementarity theorem) expresses
the factorization homology $H^\bullet(\Mbar_g, Z(\cA))$
as a direct sum of contributions from $\cA$ and its Koszul
dual $\cA^!$.
The present section lifts this additive decomposition to a
\emph{symplectic polarization}: the ambient complex carries a
natural shifted-symplectic structure, and the two summands are
Lagrangian complements.
This is the anomaly-completed interpretation of
Vol~I's ambient complementarity theorem
(Theorem~\ref*{thm:ambient-complementarity} in Vol~I).

\subsubsection{The ambient complex}
% label removed: subsubsec:tholog-ambient-complex

\begin{definition}[Ambient complex]
% label removed: def:tholog-ambient-complex
\index{ambient complex}
For a chirally Koszul algebra $\cA$ and a closed surface
$\Sigma_g$, the \emph{genus-$g$ ambient complex} is
\[
C_g(\cA)
\;:=\;
R\Gamma\!\bigl(\Mbar_g,\; Z(\cA)\bigr),
\]
where $Z(\cA)$ is the chiral center (factorization sheaf on
$\Mbar_g$).
\end{definition}

The ambient complex is a cochain complex whose cohomology is
the genus-$g$ factorization homology of $\cA$.
Its key structural feature is a natural involution.

\begin{definition}[Verdier involution on the ambient complex]
% label removed: def:tholog-verdier-involution
\index{Verdier involution!on ambient complex}
The \emph{Verdier involution} on $C_g(\cA)$ is the involution
$\sigma_g : C_g(\cA) \to C_g(\cA)$
induced by the Koszul duality $\cA \mapsto \cA^!$ at the level
of the factorization sheaf.
Explicitly, $\sigma_g$ acts on the bar complex by reversing the
bar grading: if a cochain has bar degree $p$, then
$\sigma_g$ acts by $(-1)^p$.
\end{definition}

\begin{proposition}[Eigenspace decomposition; \ClaimStatusProvedHere]
% label removed: prop:tholog-eigenspace-decomposition
The Verdier involution $\sigma_g$ is an involution
($\sigma_g^2 = \operatorname{id}$) that commutes with the
differential on $C_g(\cA)$.
The eigenspace decomposition
\[
C_g(\cA) = C_g^+(\cA) \oplus C_g^-(\cA),
\qquad
\sigma_g|_{C_g^\pm} = \pm\operatorname{id},
\]
is a direct-sum decomposition of cochain complexes.
\end{proposition}

\begin{proof}
$\sigma_g$ is an algebra involution on the bar complex, hence
commutes with the bar differential.
The factorization-homology pushforward preserves this involution
because the Koszul duality functor is compatible with the modular
operad structure (Vol~I, Theorem~A).
The eigenspace decomposition is immediate from $\sigma_g^2 = \operatorname{id}$
and $\operatorname{char}(k) = 0$.
\end{proof}

\subsubsection{The complementarity decomposition}
% label removed: subsubsec:tholog-complementarity-decomposition

\begin{theorem}[Complementarity; \ClaimStatusProvedHere]
% label removed: thm:tholog-complementarity
Let $\cA$ be a chirally Koszul algebra.
Then on the Koszul locus, there is a natural identification
\[
Q_g(\cA) \;:=\; H^\bullet(C_g^+(\cA)),
\qquad
Q_g(\cA^!) \;:=\; H^\bullet(C_g^-(\cA)),
\]
and the cohomology of the ambient complex decomposes as
\begin{equation}% label removed: eq:complementarity-decomposition
H^\bullet\!\bigl(\Mbar_g,\; Z(\cA)\bigr)
\;=\;
Q_g(\cA) \;\oplus\; Q_g(\cA^!).
\end{equation}
\end{theorem}

\begin{proof}
This is Vol~I, Theorem~C, transported to the present notation.
The $+$-eigenspace of $\sigma_g$ corresponds to the summand
computed by the bar complex of $\cA$ at degree-parity $0$,
which is identified with $Q_g(\cA)$ by the bar-cobar
quasi-isomorphism (Vol~I, Theorem~B on the Koszul locus).
The $-$-eigenspace corresponds to the bar complex at
degree-parity $1$, which is identified with $Q_g(\cA^!)$
via the Verdier intertwining of Theorem~A.
\end{proof}

\subsubsection{The cross-pairing}
% label removed: subsubsec:tholog-cross-pairing

\begin{theorem}[Cross-pairing; \ClaimStatusProvedHere]
% label removed: thm:tholog-cross-pairing
There is a natural nondegenerate pairing
\begin{equation}% label removed: eq:cross-pairing
\langle -,- \rangle_g
\;:\;
Q_g(\cA) \otimes Q_g(\cA^!)
\;\longrightarrow\;
k[-(3g-3)]
\end{equation}
of degree $-(3g-3) = -\dim_{\mathbb{C}}\Mbar_g$.
This pairing is induced by the Serre duality pairing on
$\Mbar_g$ applied to the factorization sheaf $Z(\cA)$:
\[
H^p\!\bigl(\Mbar_g, Z(\cA)\bigr)
\otimes
H^{3g-3-p}\!\bigl(\Mbar_g, Z(\cA)^\vee \otimes \omega_{\Mbar_g}\bigr)
\longrightarrow
H^{3g-3}\!\bigl(\Mbar_g, \omega_{\Mbar_g}\bigr)
\;\cong\; k.
\]
The Koszul duality $\cA \mapsto \cA^!$ identifies
$Z(\cA)^\vee$ with $Z(\cA^!)$ (up to a twist by the dualizing
sheaf correction), and the pairing descends to the
eigenspace decomposition.
\end{theorem}

\begin{proof}
The Serre duality pairing on $\Mbar_g$ (a smooth
Deligne--Mumford stack of complex dimension $3g - 3$ for
$g \geq 2$) gives a perfect pairing
\[
H^p(\Mbar_g, \cF) \otimes H^{3g-3-p}(\Mbar_g, \cF^\vee
\otimes \omega_{\Mbar_g}) \to k
\]
for any coherent sheaf $\cF$.
Taking $\cF = Z(\cA)$ and using the Koszul duality
identification $Z(\cA)^\vee \simeq Z(\cA^!)$ on the Koszul
locus (Vol~I, Theorem~A, Verdier intertwining), the pairing
\[
H^p(\Mbar_g, Z(\cA)) \otimes H^{3g-3-p}(\Mbar_g, Z(\cA^!)) \to k
\]
is well defined and nondegenerate.
Under the eigenspace decomposition
(Proposition~\ref{prop:tholog-eigenspace-decomposition}),
the $+$-eigenspace $Q_g(\cA)$ pairs nondegenerately with the
$-$-eigenspace $Q_g(\cA^!)$ (the mixed terms vanish because
$\sigma_g$ is an isometry of the Serre pairing: Serre duality
commutes with the Verdier involution, and the $\pm$-eigenspaces
are orthogonal under a pairing that intertwines with a
sign-reversing involution).

The degree shift $-(3g-3)$ accounts for the top cohomological
degree of $\Mbar_g$.
\end{proof}

\subsubsection{The shifted-symplectic structure}
% label removed: subsubsec:tholog-shifted-symplectic

\begin{theorem}[Shifted-symplectic structure on the ambient complex; \ClaimStatusProvedHere]
% label removed: thm:tholog-shifted-symplectic
\index{shifted-symplectic structure}
Assume $\cA$ is chirally Koszul with nondegenerate cyclic
pairing.
The ambient complex $C_g(\cA)$ carries a
$(-1)$-shifted symplectic structure $\omega_{\mathrm{amb}}$:
a closed degree-$(-1)$ bilinear form
\[
\omega_{\mathrm{amb}}
\;:\;
C_g(\cA) \otimes C_g(\cA) \to k[-(3g-2)]
\]
that is nondegenerate on cohomology.
Explicitly, $\omega_{\mathrm{amb}}$ is the composition
\begin{equation}% label removed: eq:symplectic-form
\omega_{\mathrm{amb}}(x, y)
\;=\;
\langle x,\; \sigma_g(y) \rangle_{\mathrm{Serre}}
\;-\;
(-1)^{|x||y|}\langle y,\; \sigma_g(x) \rangle_{\mathrm{Serre}},
\end{equation}
where $\langle -, - \rangle_{\mathrm{Serre}}$ is the Serre
duality pairing of Theorem~\ref{thm:tholog-cross-pairing}.
\end{theorem}

\begin{proof}
The form $\omega_{\mathrm{amb}}$ is manifestly antisymmetric
(in the graded sense).
We verify nondegeneracy and closedness.

\emph{Nondegeneracy.}
Decompose $x = x^+ + x^-$ and $y = y^+ + y^-$ into
$\sigma_g$-eigenspaces.
Then $\sigma_g(y) = y^+ - y^-$, and
\begin{align*}
\omega_{\mathrm{amb}}(x, y)
&= \langle x^+ + x^-,\; y^+ - y^- \rangle_{\mathrm{Serre}}
- (-1)^{|x||y|}\langle y^+ + y^-,\; x^+ - x^- \rangle_{\mathrm{Serre}} \\
&= \langle x^+, -y^- \rangle + \langle x^-, y^+ \rangle
- (-1)^{|x||y|}(\langle y^+, -x^- \rangle
+ \langle y^-, x^+ \rangle),
\end{align*}
where we used that the Serre pairing vanishes on same-eigenspace
pairs ($\langle x^+, y^+ \rangle = 0$ and
$\langle x^-, y^- \rangle = 0$, because $\sigma_g$ acts by the
same sign on both entries, but the pairing satisfies
$\langle \sigma_g x, \sigma_g y \rangle = \langle x, y \rangle$
and antisymmetry of $\omega_{\mathrm{amb}}$ forces the
same-eigenspace part to vanish).

The cross-pairing terms $\langle x^-, y^+ \rangle$ and
$\langle x^+, y^- \rangle$ are nondegenerate by
Theorem~\ref{thm:tholog-cross-pairing}, hence
$\omega_{\mathrm{amb}}$ is nondegenerate.

\emph{Closedness.}
The form $\omega_{\mathrm{amb}}$ is built from the
Serre pairing (which is $d$-closed on cohomology) and
the Verdier involution $\sigma_g$ (which commutes with $d$).
Therefore $d\omega_{\mathrm{amb}} = 0$ on the cochain level.
\end{proof}

\subsubsection{Lagrangian condition}
% label removed: subsubsec:tholog-lagrangian

\begin{theorem}[Lagrangian complementarity; \ClaimStatusProvedHere]
% label removed: thm:tholog-lagrangian
\index{Lagrangian!complementarity!proof}
The eigenspaces $C_g^+(\cA)$ and $C_g^-(\cA)$ are
Lagrangian subcomplexes of $(C_g(\cA), \omega_{\mathrm{amb}})$:
\begin{enumerate}[label=\textup{(\roman*)}]
\item $\omega_{\mathrm{amb}}|_{C_g^+ \otimes C_g^+} = 0$.
\item $\omega_{\mathrm{amb}}|_{C_g^- \otimes C_g^-} = 0$.
\item $C_g^+ \oplus C_g^- = C_g(\cA)$.
\item The restricted pairing
  $\omega_{\mathrm{amb}}|_{C_g^+ \otimes C_g^-}$
  is nondegenerate.
\end{enumerate}
In the language of derived algebraic geometry:
$(Q_g(\cA), Q_g(\cA^!))$ is a Lagrangian pair in the
$(-1)$-shifted symplectic space
$H^\bullet(\Mbar_g, Z(\cA))$.
\end{theorem}

\begin{proof}
Parts~(i) and~(ii): For $x, y$ in the same eigenspace,
$\sigma_g(y) = \pm y$, and the computation in the proof of
Theorem~\ref{thm:tholog-shifted-symplectic} shows that
$\omega_{\mathrm{amb}}(x, y) = 0$
when both $x$ and $y$ belong to the same eigenspace.
This is because the Serre pairing
$\langle x^\pm, y^\pm \rangle = 0$ vanishes on same-eigenspace pairs
(the Verdier involution acts as $\sigma_g(y^\pm) = \pm y^\pm$, and
$\omega_{\mathrm{amb}}(x^\pm, y^\pm)
= \langle x^\pm, \pm y^\pm \rangle
\mp (-1)^{|x||y|}\langle y^\pm, \pm x^\pm \rangle
= \pm(\langle x^\pm, y^\pm \rangle
- (-1)^{|x||y|}\langle y^\pm, x^\pm \rangle)$.
By graded symmetry of the Serre pairing:
$\langle y^\pm, x^\pm \rangle
= (-1)^{|x||y|}\langle x^\pm, y^\pm \rangle$,
so the expression is
$\pm(\langle x^\pm, y^\pm \rangle
- \langle x^\pm, y^\pm \rangle) = 0$.)

Part~(iii) is Proposition~\ref{prop:tholog-eigenspace-decomposition}.

Part~(iv) follows from the nondegeneracy of
$\omega_{\mathrm{amb}}$ (Theorem~\ref{thm:tholog-shifted-symplectic})
together with the isotropy of (i)--(ii): the only way a
nondegenerate form can vanish on two complementary subspaces
is if the cross-pairing is nondegenerate.
\end{proof}

\begin{remark}[Anti-symplecticity of the Verdier involution]
% label removed: rem:tholog-anti-symplectic
The Verdier involution $\sigma_g$ is anti-symplectic:
$\omega_{\mathrm{amb}}(\sigma_g x, \sigma_g y)
= -\omega_{\mathrm{amb}}(x, y)$.
This follows from the explicit
formula~\eqref{eq:symplectic-form}:
replacing $(x, y)$ by $(\sigma_g x, \sigma_g y)$ flips
the sign because $\sigma_g^2 = \operatorname{id}$ and the
form is defined by $\langle x, \sigma_g y \rangle$ (which
becomes $\langle \sigma_g x, y \rangle$, swapping the
roles of $Q_g(\cA)$ and $Q_g(\cA^!)$).
The anti-symplecticity is the structural reason the
eigenspaces are Lagrangian: any involution anti-symplectic
with respect to a nondegenerate $2$-form necessarily has
isotropic fixed-point and anti-fixed-point loci.
\end{remark}

\subsubsection{Genus 1 for Virasoro: explicit dimensions}
% label removed: subsubsec:tholog-genus1-virasoro

\begin{computation}[Genus-1 complementarity for Virasoro; \ClaimStatusProvedHere]
% label removed: comp:tholog-virasoro-genus1
\index{Virasoro algebra!genus-1 complementarity}
At genus $g = 1$, $\Mbar_1 \cong \mathbb{P}^1$
(after adding the nodal curve), $\dim_\mathbb{C} = 1$.
For $\mathrm{Vir}_c$ (generic $c$, away from $c = 0, 13, 26$):

\emph{The ambient complex.}
$C_1(\mathrm{Vir}_c) = R\Gamma(\Mbar_1, Z(\mathrm{Vir}_c))$.
The factorization sheaf $Z(\mathrm{Vir}_c)$ on $\Mbar_1$
is determined by the Virasoro bar complex restricted to
genus-$1$ data: the one-loop amplitudes.
The relevant cohomology is
\[
H^\bullet(\Mbar_1, Z(\mathrm{Vir}_c))
\;=\;
Q_1(\mathrm{Vir}_c) \;\oplus\; Q_1(\mathrm{Vir}_{26-c}).
\]

\emph{Dimensions.}
The genus-$1$ partition function of $\mathrm{Vir}_c$ is the
generating function of the Virasoro characters at genus~$1$.
The relevant invariants:
\begin{enumerate}[label=\textup{(\roman*)}]
\item $\dim Q_1(\mathrm{Vir}_c)$ counts the genus-$1$
  Virasoro modular data at central charge $c$:
  the image of the torus partition function in the
  quotient by the Koszul dual contribution.
\item $\dim Q_1(\mathrm{Vir}_{26-c})$ counts the dual
  contribution at the complementary charge.
\item The total dimension satisfies
  \[
  \dim Q_1(\mathrm{Vir}_c) + \dim Q_1(\mathrm{Vir}_{26-c})
  = \dim H^\bullet(\Mbar_1, Z(\mathrm{Vir}_c)).
  \]
\item At the self-dual point $c = 13$:
  $Q_1(\mathrm{Vir}_{13}) \cong Q_1(\mathrm{Vir}_{13})$
  (isomorphic, by the self-duality),
  so each Lagrangian has exactly half the total dimension.
\item At $c = 26$: the Koszul dual $\mathrm{Vir}_0$ has
  $T = 0$, so $Q_1(\mathrm{Vir}_0) = 0$ and
  $Q_1(\mathrm{Vir}_{26})$ carries the full ambient cohomology.
  The complementarity decomposition degenerates: one Lagrangian
  is the full space, the other is zero.
\end{enumerate}

\emph{The cross-pairing.}
At genus $1$, $\dim_\mathbb{C}\Mbar_1 = 1$, so the
cross-pairing has degree $-(3 \cdot 1 - 3) = 0$:
\[
\langle -, - \rangle_1
\;:\;
Q_1(\mathrm{Vir}_c) \otimes Q_1(\mathrm{Vir}_{26-c})
\;\longrightarrow\; k.
\]
This is a degree-$0$ pairing.
Its explicit evaluation on the genus-$1$ modular form
$\chi_c(\tau) = q^{-c/24}\prod_{n=1}^\infty(1-q^n)^{-1}$:
the cross-pairing at genus $1$ is controlled by the
Petersson inner product on the space of weakly holomorphic
modular forms, restricted to the Virasoro characters
at complementary central charges.
\end{computation}

\subsubsection{The complementarity potential}
% label removed: subsubsec:tholog-complementarity-potential

\begin{definition}[Complementarity potential]
% label removed: def:tholog-complementarity-potential
\index{complementarity potential}
Suppose the Lagrangian $Q_g(\cA^!)$ can be written as the
graph of a $1$-form on $Q_g(\cA)$:
there exists a functional
\[
S_\cA \;:\; Q_g(\cA) \to k[-(3g-3)]
\]
(the \emph{complementarity potential}) such that
\[
Q_g(\cA^!) = \bigl\{(x, dS_\cA(x)) \;\big|\; x \in Q_g(\cA)\bigr\}
\]
under the symplectic identification
$C_g(\cA) = Q_g(\cA) \oplus Q_g(\cA)^\vee[-(3g-3)]$.
\end{definition}

\begin{proposition}[Existence and uniqueness; \ClaimStatusProvedHere]
% label removed: prop:tholog-potential-existence
On the Koszul locus, the complementarity potential $S_\cA$
exists and is unique up to constants.
Its existence is equivalent to the transversality of the two
Lagrangians: $Q_g(\cA) \cap Q_g(\cA^!) = \{0\}$.
\end{proposition}

\begin{proof}
The symplectic vector space $C_g(\cA) = L_+ \oplus L_-$
with $L_\pm$ Lagrangian and $L_+ \cap L_- = \{0\}$ admits
the identification $C_g(\cA) \cong L_+ \oplus L_+^\vee$
via the symplectic form.
The second Lagrangian $L_-$, being transverse to $L_+$,
is the graph of a symmetric bilinear form on $L_+$
(i.e.\ a quadratic function $S_\cA$ on $L_+$, up to a constant).
The form is symmetric because $L_-$ is Lagrangian (isotropy
translates to the symmetry of the quadratic form).
Uniqueness up to constants is standard.
\end{proof}

\subsubsection{Shadow jets as Taylor coefficients}
% label removed: subsubsec:tholog-shadow-jets

\begin{theorem}[Shadow jets from the complementarity potential; \ClaimStatusProvedHere]
% label removed: thm:tholog-shadow-jets
The Taylor expansion of the complementarity potential $S_\cA$
around the origin produces the shadow jets of Vol~I:
\begin{equation}% label removed: eq:shadow-jet-expansion
S_\cA(x) = \sum_{r=2}^{\infty} \frac{1}{r!}\,
\mathrm{Sh}_r(\Theta_\cA)(x^{\otimes r}),
\end{equation}
where $\mathrm{Sh}_r(\Theta_\cA)$ is the arity-$r$ shadow
projection of the MC element $\Theta_\cA \in \mathrm{MC}(g_\cA^{\mathrm{mod}})$.
In particular:
\begin{enumerate}[label=\textup{(\roman*)}]
\item The quadratic term
  $\frac{1}{2}\mathrm{Sh}_2(\Theta_\cA)(x, x)
  = \kappa(\cA) \cdot \langle x, x \rangle$
  is the modular Koszul curvature.
\item The cubic term
  $\frac{1}{6}\mathrm{Sh}_3(\Theta_\cA)(x, x, x)
  = C(\cA)(x^{\otimes 3})$
  is the cubic shadow.
\item The quartic term
  $\frac{1}{24}\mathrm{Sh}_4(\Theta_\cA)(x^{\otimes 4})
  = Q(\cA)(x^{\otimes 4})$
  is the quartic resonance class.
\end{enumerate}
The expansion~\eqref{eq:shadow-jet-expansion} converges
(in the sense of the arity filtration) to the full
MC element $\Theta_\cA$ by Vol~I,
Theorem~\textup{\ref*{V1-thm:recursive-existence}}
(all-arity convergence).
\end{theorem}

\begin{proof}
The complementarity potential $S_\cA$ is determined by the
condition that $Q_g(\cA^!)$ is the graph of $dS_\cA$.
The dual Lagrangian $Q_g(\cA^!)$ is in turn determined by
the bar-cobar inversion applied to the MC element $\Theta_\cA$.

The Taylor expansion of $S_\cA$ at the origin is computed by
expanding the MC equation
$D \Theta_\cA + \frac{1}{2}[\Theta_\cA, \Theta_\cA] = 0$
in the arity filtration.
At arity $r$, the MC equation projected to degree $r$ involves
only the shadow projections
$\mathrm{Sh}_2, \dots, \mathrm{Sh}_r$, and the arity-$r$
Taylor coefficient of $S_\cA$ is
$\frac{1}{r!}\mathrm{Sh}_r(\Theta_\cA)$.

Part~(i): at arity $2$, the MC equation reduces to the
curvature equation $d^2 = \kappa \cdot \omega_g$
(Vol~I, Theorem~D), and the quadratic Taylor coefficient
of $S_\cA$ is the scalar curvature $\kappa(\cA)$.

Part~(ii): at arity $3$, the MC equation gives the cubic
shadow $C(\cA)$ as the obstruction class to trivializing the
arity-$3$ truncation (Vol~I, cubic gauge triviality,
Theorem~\ref*{thm:cubic-gauge-triviality}).

Part~(iii): at arity $4$, the MC equation gives the quartic
resonance class $Q(\cA)$ (Vol~I, quartic resonance with
clutching law via Mok's log FM degeneration).

Convergence follows from Vol~I,
Theorem~\ref*{V1-thm:recursive-existence}:
$\Theta_\cA = \varprojlim \Theta_\cA^{\leq r}$.
\end{proof}

\begin{remark}[The complementarity potential as a generating function]
% label removed: rem:tholog-potential-generating-function
The complementarity potential $S_\cA$ is the generating
function of all shadow invariants simultaneously.
Its content is equivalent to the full MC element $\Theta_\cA$,
but organized differently: $\Theta_\cA$ lives in the
convolution dg Lie algebra $g_\cA^{\mathrm{mod}}$, while
$S_\cA$ lives on the Lagrangian complement $Q_g(\cA)$.
The translation between the two is the Legendre transform
of the MC equation, passing from the Lie-algebraic
formulation (``$\Theta$ satisfies the MC equation'') to
the variational formulation (``$Q_g(\cA^!)$ is the critical
locus of $S_\cA + \omega_{\mathrm{amb}}$'').
This is the Chriss--Ginzburg principle at the level of
generating functions: every algebraic structure in the
monograph is an MC element, and the complementarity potential
is the Lagrangian shadow of that MC element.
\end{remark}

\subsubsection{Perfectness and nondegeneracy}
% label removed: subsubsec:tholog-perfectness

\begin{definition}[Perfect ambient complex]
% label removed: def:tholog-perfect-ambient
\index{perfect ambient complex}
The ambient complex $C_g(\cA)$ is \emph{perfect} if it is
quasi-isomorphic to a bounded complex of finitely generated
free $k$-modules.
It is \emph{nondegenerate} if in addition the cross-pairing
$\langle -, - \rangle_g$ of
Theorem~\ref{thm:tholog-cross-pairing}
is a perfect pairing (i.e.\ induces an isomorphism
$Q_g(\cA) \xrightarrow{\sim} Q_g(\cA^!)^\vee[-(3g-3)]$).
\end{definition}

\begin{proposition}[Perfectness for the standard landscape; \ClaimStatusProvedHere]
% label removed: prop:tholog-perfectness-landscape
The ambient complex $C_g(\cA)$ is perfect for $\cA$ in the
standard landscape (Heisenberg, affine Kac--Moody at
non-critical level, $\beta\gamma$, Virasoro at generic $c$,
$\mathcal{W}_N$ at generic $c$) and for all $g \geq 0$.
\end{proposition}

\begin{proof}
Perfectness follows from the finiteness properties of the
bar complex.
For each standard family, the bar complex is the cochain complex
of a finitely generated operadic algebra with polynomial growth
of the PBW basis (Vol~I, MC1: PBW concentration for all standard
families).
The factorization-homology pushforward to $\Mbar_g$ preserves
perfectness because $\Mbar_g$ is a smooth proper
Deligne--Mumford stack (Grothendieck duality for proper morphisms).
\end{proof}

\begin{theorem}[Nondegeneracy implies shifted-symplectic; \ClaimStatusProvedHere]
% label removed: thm:tholog-nondegeneracy-symplectic
If $C_g(\cA)$ is both perfect and nondegenerate, then:
\begin{enumerate}[label=\textup{(\roman*)}]
\item The $(-1)$-shifted symplectic form $\omega_{\mathrm{amb}}$
  of Theorem~\ref{thm:tholog-shifted-symplectic} is
  nondegenerate at the cochain level (not merely on cohomology).
\item The Lagrangian pair $(Q_g(\cA), Q_g(\cA^!))$ extends
  to a Lagrangian pair at the derived level:
  $(C_g^+(\cA), C_g^-(\cA))$ are derived Lagrangian
  subcomplexes.
\item The complementarity potential $S_\cA$ extends to a
  cochain-level functional
  $S_\cA^{\mathrm{ch}} : C_g^+(\cA) \to k[-(3g-2)]$
  satisfying $dS_\cA^{\mathrm{ch}} = 0$ and inducing
  $S_\cA$ on cohomology.
\end{enumerate}
\end{theorem}

\begin{proof}
Part~(i): by the nondegeneracy hypothesis, the cross-pairing
$Q_g(\cA) \otimes Q_g(\cA^!) \to k[-(3g-3)]$ is perfect.
Since $C_g(\cA)$ is perfect, the symplectic form is
represented by a cochain-level pairing that is nondegenerate
on a cofibrant replacement.

Part~(ii): the eigenspaces $C_g^\pm(\cA)$ are subcomplexes
with the isotropy property proved in
Theorem~\ref{thm:tholog-lagrangian}, and the nondegeneracy
of the ambient form restricted to the cross-terms
$C_g^+ \otimes C_g^-$ holds at the cochain level by
perfectness.

Part~(iii): the complementarity potential
$S_\cA$ extends to the cochain level by the same graph
construction as Proposition~\ref{prop:tholog-potential-existence},
applied to the cochain-level Lagrangians.
Closedness $dS_\cA^{\mathrm{ch}} = 0$ follows from the
Lagrangian condition (isotropy at the cochain level).
\end{proof}

\subsubsection{The symplectic polarization as a functor}
% label removed: subsubsec:tholog-symplectic-functor

\begin{proposition}[Functoriality; \ClaimStatusProvedHere]
% label removed: prop:tholog-complementarity-functorial
The assignment
\[
\cA \longmapsto
\bigl(C_g(\cA),\; \omega_{\mathrm{amb}},\;
Q_g(\cA),\; Q_g(\cA^!)\bigr)
\]
is functorial in the following sense:
for any morphism $f : \cA \to \cB$ of chirally Koszul
algebras (compatible with the Koszul duality), there is an
induced symplectomorphism
\[
f_* : C_g(\cA) \to C_g(\cB)
\]
preserving the Lagrangian decomposition.
If $f$ is a quasi-isomorphism, then $f_*$ is a
quasi-isomorphism of $(-1)$-shifted symplectic complexes.
\end{proposition}

\begin{proof}
The factorization-homology pushforward is functorial in the
input algebra, and the Verdier involution is natural.
The symplectic form is defined in terms of the Serre pairing
(which is natural) and the Verdier involution (which is
natural by Theorem~A).
The Lagrangian decomposition is the eigenspace decomposition
under $\sigma_g$, which is preserved by $f_*$ because
$\sigma_g$ is natural.
\end{proof}

\begin{remark}[Summary of the symplectic polarization]
% label removed: rem:tholog-symplectic-summary
\index{complementarity!summary}
The full picture:
\begin{enumerate}[label=\textup{(\arabic*)}]
\item The ambient complex
  $C_g(\cA) = R\Gamma(\Mbar_g, Z(\cA))$
  carries a $(-1)$-shifted symplectic structure
  from Serre duality on $\Mbar_g$.
\item The Verdier involution $\sigma_g$ (from Koszul duality)
  is anti-symplectic, producing a Lagrangian pair
  $(Q_g(\cA), Q_g(\cA^!))$.
\item The cross-pairing
  $Q_g(\cA) \otimes Q_g(\cA^!) \to k[-(3g-3)]$
  is perfect on the Koszul locus.
\item The complementarity potential $S_\cA$ is the generating
  function whose Taylor coefficients are the shadow jets
  $\kappa, C, Q, \dots$
\item At the critical point ($\kappa = 0$), the Lagrangian
  pair degenerates: $Q_g(\cA) \cong Q_g(\cA^!)$ and
  $S_\cA = 0$.  (For KM this is the self-dual point;
  for $\mathcal{W}$-algebras the self-dual point has
  $\kappa = \kappa^! \neq 0$.)
\item On the gravitational locus ($u = 0$), the genus-Clifford
  completion degenerates to an exterior algebra, and the
  symplectic polarization reduces to the standard Hodge
  decomposition on $\Lambda^\bullet(H_1(\Sigma_g))$.
\end{enumerate}
\end{remark}
